\documentclass[12pt,a4paper,oneside]{report}

\usepackage[T1]{fontenc}
\usepackage[a4paper,margin=2.5cm]{geometry}
\usepackage{setspace}
\onehalfspacing
\usepackage{microtype}
\usepackage{amsmath,amssymb,mathtools}
\usepackage{booktabs,tabularx,longtable,array,multirow}
\usepackage{enumitem}
\usepackage{xcolor}
\usepackage[most]{tcolorbox}
\usepackage{csquotes}
\usepackage{graphicx}
\usepackage{caption}
\usepackage{subcaption}
\usepackage{hyperref}
\usepackage[nameinlink,noabbrev]{cleveref}
\usepackage[backend=biber,style=authoryear,sorting=nyt,maxcitenames=2,maxbibnames=99,giveninits=true,uniquename=false,doi=true,url=true,isbn=false]{biblatex}
\addbibresource{selective_communication_references.bib}

\hypersetup{
  colorlinks=true,
  linkcolor=blue!45!black,
  citecolor=blue!45!black,
  urlcolor=blue!45!black,
  pdftitle={Information Constraints and Institutional Objectives in Financial-Assistant LLM Communication}
}

\newcolumntype{Y}{>{\raggedright\arraybackslash}X}
\newcolumntype{L}[1]{>{\raggedright\arraybackslash}p{#1}}

\newtcolorbox{reviewnote}[1][]{
  breakable,
  enhanced,
  colback=yellow!18,
  colframe=yellow!45!black,
  fonttitle=\bfseries,
  title={Review note},
  boxrule=0.7pt,
  arc=1mm,
  left=2mm,right=2mm,top=1.5mm,bottom=1.5mm,
  #1
}

\newcommand{\code}[1]{\texttt{#1}}
\newcommand{\owner}{deploying financial institution}
\newcommand{\RQ}[1]{\par\noindent\textbf{#1}\quad}

\title{Information Constraints and Institutional Objectives in Financial-Assistant LLM Communication}
\author{\textbf{Iman Zafar}}
\date{Master's Dissertation Draft\\9 August 2026}

\begin{document}

\pagenumbering{roman}
\maketitle

\tableofcontents
\clearpage
\pagenumbering{arabic}

\chapter{Introduction}
\label{chap:introduction}

Financial services is an information-intensive sector in which customers are routinely required to compare costs, conditions, risks and benefits across products or courses of action. Large language models (LLMs) are increasingly capable of mediating these decisions through conversational interfaces. Their potential role extends beyond answering isolated questions. A financial assistant can identify relevant information, compare available options, explain technical terms, respond to follow-up questions and help a customer determine what requires further attention. These capabilities make conversational systems attractive as a scalable interface between complex financial information and customers who may have neither the time nor the technical knowledge to interpret the underlying documentation unaided \parencite{li_etal_2023_finance_survey}.

The prospect of greater use of such systems is no longer confined to experimental settings. The Bank of England and Financial Conduct Authority reported widespread adoption of artificial intelligence across UK financial services in 2024, with further growth expected across a range of functions including customer support \parencite{boe_fca_2024_ai}. More recently, the Mills Review identified consumer-facing AI as a potentially important part of future retail financial journeys, particularly where choices are complex or consequential, and noted that AI systems may increasingly provide comparisons, recommendations and eventually act on consumers' behalf \parencite{fca_2026_mills_review}. This creates an important evaluation problem. As conversational systems become more involved in financial decisions, their quality cannot be assessed solely by asking whether the individual statements they produce are factually correct.

A customer-facing assistant also determines which available facts enter an answer, how much detail is retained, which considerations appear first and how limited response space is allocated between competing considerations. These choices matter because communication can be incomplete or directionally unbalanced while remaining literally truthful. A response might correctly describe a lower interest rate while omitting a substantial product fee, mention that funds are accessible while failing to state the notice period, or discuss the advantages of one option in considerably more detail than its disadvantages. In each case, the issue lies not in the truth value of the statements that appear, but in the relationship between the information available to the model and the subset of that information ultimately presented to the customer.

This distinction has particular importance in financial services. The Financial Conduct Authority's Consumer Duty requires firms to support customers in making effective and properly informed decisions, including through communications that meet their information needs \parencite{fca_2022_consumer_duty}. The FCA has maintained that existing outcomes-based requirements continue to apply when firms use artificial intelligence \parencite{fca_2026_ai_approach}. From this perspective, an answer may create a meaningful communication risk even when it contains no hallucinated figure or directly false statement. A material qualification can be omitted, an otherwise accurate fact can lose the numerical detail that makes it useful, or a comparison can become one-sided because limited space is allocated unevenly.

The possibility of truthful but selectively incomplete communication also connects this problem to a broader literature on deception and information manipulation. Research on AI deception has shown that language models can conceal information or misrepresent their behaviour in settings where doing so advances an objective \parencite{park_etal_2024_ai_deception,scheurer_etal_2024_strategic_deception}. More recent work has shifted attention from direct falsehood towards manipulation of otherwise available information. JANUS examines omission, weakening and differential preservation of favourable and adverse facts under goal-conditioned generation, while DECOR decomposes responses into informational units to identify more granular forms of distortion \parencite{giannouris_etal_2026_janus,cai_etal_2026_decor}. These approaches are important because they recognise that a misleading overall representation need not depend on fabrication.

The present study focuses on a narrower behavioural question. It does not attempt to infer whether an LLM intends to deceive a customer. Instead, it asks how materially relevant financial information is selected and presented when the conditions under which the model communicates change. Two distinctions are central. First, \emph{information imbalance} is separated from \emph{direction}. A constrained answer may become incomplete or one-sided without systematically favouring the institution represented by the assistant. Second, an institutional affiliation is separated from an explicit institutional objective. A model assigned to act as a support assistant for a fictional financial institution may not behave in the same way as a model that is directly instructed to protect that institution's commercial interests.

These distinctions motivate three groups of experimental conditions. The first concerns \emph{customer state}. Financial questions may be expressed neutrally or may contain signs of anxiety or frustration, and conversational models are capable of adapting their responses to social and affective cues \parencite{wu_etal_2024_social_signal,hartley_etal_2025_risk}. Such adaptation might alter only tone and reassurance, or it might also change which material considerations are prioritised. The second concerns \emph{information constraints}. Real conversational interfaces frequently favour short responses, yet reducing available space requires information to be prioritised. The study therefore distinguishes exact limits on the number of facts that may be selected from more natural word-count constraints and from tasks that ask the model to identify a single priority. The third concerns \emph{institutional objectives}. The study separately examines affiliation with a fictional provider and an explicit instruction to protect that provider's commercial interests.

To study these questions under controlled conditions, a synthetic financial benchmark was constructed containing 30 scenarios across six domains: mortgages, credit and repayment, savings, investment platforms, insurance settlements and international payments. Every scenario presents two mutually exclusive options and six material facts organised into three matched pairs. Each pair contains one fact relating to each option and matches the two facts on customer valence. Across the corpus, institutional direction is also balanced: three facts in each scenario support the option that provides greater benefit to the deploying institution and three countervail that option. The evaluated model does not receive these research labels. It sees named fictional institutions, the available options, the six factual statements and a natural customer query.

Seven language models are evaluated across seven experiments. These experiments manipulate customer-state wording, exact information budgets, natural response-length limits, first-priority requests, institutional role, forced option choice and an explicit commercial-interest instruction. The study therefore does not treat concision as a single binary intervention. Instead, it examines whether different forms of restricted communication produce general information loss, pairwise imbalance, a systematic directional shift, or some combination of these outcomes. Exact-budget tasks additionally require models to identify the facts they have selected before producing prose, allowing prioritisation itself to be distinguished from the subsequent realisation of that selection in a customer-facing answer.

The measurement framework follows the same separation of constructs. The principal direction-sensitive outcome, \(D\), measures the net difference between owner-supporting and countervailing facts communicated across the three matched pairs. Absolute pairwise imbalance, \(A\), records whether matched facts are communicated unevenly regardless of direction, while total material coverage, \(T\), records how much of the available fact set survives into the response. Specificity, factual presentation, recommendation behaviour and factual errors are analysed separately rather than being combined into a single index. This distinction is important because a shorter response may become substantially less complete without becoming institution-favouring, while an institutional objective may create a directional effect even where overall coverage changes little.

The resulting design makes three contributions. First, it provides a controlled financial communication benchmark in which customer valence, institutional direction and matched informational trade-offs are represented separately. Second, it distinguishes general selectivity caused by limited communication space from selection that systematically favours an institutional interest. Third, across seven language models, the experiments provide evidence that institutional affiliation and an explicit commercial objective affect communication differently: affiliation alone primarily changes ordering, whereas an explicit commercial-interest instruction changes framing and recommendation and, under tight information constraints, the direction of factual selection. Together, these features provide a framework for examining how ordinary communication constraints and institutional objectives shape financial-assistant responses without treating every omission as deception or every institutional association as evidence of an internal commercial goal.

\section{Research questions}
\label{sec:research_questions}

The preceding discussion motivates three research questions concerning the conditions under which financial-assistant LLMs alter the selection and presentation of materially relevant information. The questions distinguish between three conceptually separate influences on communication: the institutional context within which the assistant is instructed to operate, the state expressed by the customer, and constraints on the amount of information that can be conveyed. Across these questions, the analysis separates directional information selection from general imbalance, overall coverage, specificity, presentation and factual accuracy.

\RQ{RQ1.} How do institutional affiliation and an explicit commercial-interest objective affect the selection and presentation of material information by financial-assistant LLMs?

This question examines whether institutional alignment emerges from the assistant's assigned role, from an explicit statement of commercial interest, or from the interaction between institutional objectives and tasks requiring information prioritisation. The commercial-interest experiment compares otherwise matched prompts with and without an instruction to protect the employing institution's commercial interests. The institutional-role experiment separately varies whether the assistant is associated with either of two fictional providers or positioned as an independent adviser. The analysis considers whether these conditions alter directional fact selection, pairwise imbalance, factual coverage, specificity, presentation, recommendation behaviour and factual accuracy. Institutional affiliation is part of the scenario context and is held constant across matched conditions unless it is itself the factor varied by the ownership-role or ownership-flip design.

\RQ{RQ2.} How do customer-state cues affect the selection and presentation of materially relevant information by financial-assistant LLMs?

This question examines whether otherwise comparable financial queries expressed in neutral, anxious or frustrated terms produce systematic differences in the information communicated by the model. The principal directional comparison concerns whether anxious wording changes the relative communication of owner-supporting and countervailing facts compared with neutral wording. The wider analysis considers whether customer state also affects pairwise information balance, overall material coverage, specificity, response presentation, reassurance, referral behaviour and factual accuracy. Query length is varied independently so that adaptation to customer state can be distinguished from differences associated with the amount of contextual language in the customer's request.

\RQ{RQ3.} How do constraints on response space and explicit prioritisation requirements affect the balance, direction and completeness of material information communicated by financial-assistant LLMs?

This question examines whether limiting the information that can be selected or expressed changes not only how much material survives into a response, but also which information is prioritised. Exact information budgets provide a controlled measure of selection when models are permitted to choose (k=6), (k=4) or (k=2) facts from the available set. Natural 40-, 80- and 160-word response limits provide a complementary test under more conventional conversational constraints, while single-priority and option-selection tasks examine behaviour when the model is required to identify one consideration or one option as most salient. These experiments allow general information loss and pairwise imbalance to be distinguished from a systematic directional preference for owner-supporting information.

Taken together, the three research questions separate effects arising from the institution represented by the assistant, the customer and the communication constraint. This structure is important because selective communication need not be institutionally directional. A model may omit substantial information when response space is restricted while remaining balanced between owner-supporting and countervailing facts. Conversely, an institutional objective may alter the direction in which information is selected even where the overall amount of material communicated changes little. The experimental design therefore evaluates these dimensions separately rather than treating all forms of omission, prioritisation or presentation asymmetry as a single construct.

\chapter{Literature Review}
\label{chap:literature_review}

\section{Customer-facing language models in financial services}

Large language models have extended the range of tasks for which natural-language systems can be used in finance. Existing research covers financial question answering, document analysis, summarisation, information extraction, forecasting support and decision-support applications, alongside persistent concerns about domain knowledge, reliability and governance \parencite{li_etal_2023_finance_survey}. Customer-facing use introduces an additional problem because the output is not simply an internal analytic artefact. It becomes part of the information environment through which a consumer interprets a financial decision.

This role is likely to become more important as artificial intelligence is incorporated more deeply into financial services. The 2024 Bank of England and FCA survey found widespread use of AI among responding firms and expectations of further expansion across the sector \parencite{boe_fca_2024_ai}. The Mills Review subsequently described a possible shift towards AI-mediated and increasingly delegated retail financial journeys, in which systems may compare products, make recommendations and, over time, undertake actions on consumers' behalf \parencite{fca_2026_mills_review}. Although these developments do not imply that current general-purpose LLMs are already operating as autonomous financial advisers, they establish the relevance of studying how conversational models organise information around consequential choices.

The regulatory concern is similarly broader than factual correctness. Under the FCA's Consumer Duty, firms are expected to support retail customers in making effective, timely and properly informed decisions \parencite{fca_2022_consumer_duty}. The FCA's technology-neutral approach means that the use of AI does not displace these existing obligations \parencite{fca_2026_ai_approach}. A customer communication may therefore be deficient without containing an obvious falsehood. It can state accurate facts but omit a material qualification, present one option with substantially greater precision, or allocate most of the available explanation to one side of a decision.

This motivates an evaluation framework in which the model's \emph{selection} of information is itself an outcome. Conventional task benchmarks generally define success through a correct answer, a retrieved fact, a numerical calculation or agreement with a labelled target. Those measures remain important, but they do not reveal whether an assistant supplied with several decision-relevant facts treats them symmetrically when constructing a response. In financial communication, the distinction is consequential because users may reasonably interpret inclusion, omission and prominence as cues about what deserves attention.

\section{Selective communication, deception and information manipulation}

The idea that communication can mislead without depending on a direct falsehood predates contemporary work on language models. Grice's account of cooperative conversation distinguishes expectations concerning the quantity, quality, relevance and manner of information supplied in an exchange \parencite{grice_1975_logic_conversation}. Building on this pragmatic tradition, Information Manipulation Theory argues that misleading communication can operate through manipulation of the amount, truthfulness, relevance or clarity of available information \parencite{mccornack_1992_information_manipulation}. This literature is useful for the present problem because it treats omission and selective presentation as features of communication in their own right rather than as imperfect versions of lying.

Research on artificial intelligence has often adopted a narrower conception of deception involving strategically false or concealed information. \textcite{park_etal_2024_ai_deception} survey deceptive behaviour across AI systems and distinguish it from ordinary model error or hallucination. Empirical work has shown that language models can conceal actions or misrepresent information when deceptive behaviour is useful for an assigned objective. In the simulated financial environment studied by \textcite{scheurer_etal_2024_strategic_deception}, for example, an autonomous trading agent engaged in prohibited conduct under organisational pressure and subsequently gave a misleading explanation of its actions. Related work has examined how reinforced agents learn strategically deceptive communication in simulated environments \parencite{dogra_etal_2024_reinforced_deception}.

These studies demonstrate an important capability, but their experimental settings usually make the conflict or objective comparatively explicit. That is appropriate when the research question concerns strategic deception, but it does not cover the full range of communication risks that can arise in ordinary support interactions. A customer-facing assistant may simply be asked to compare two products from a supplied information set. If one material downside disappears from its answer, the omission may affect the usefulness of the response without providing evidence that the model formed an intention to deceive.

Recent LLM benchmarks have therefore moved towards more granular accounts of information distortion. JANUS presents models with fixed sets of favourable and adverse facts and examines how goal conditioning affects omission, weakening, emphasis and retention of source information \parencite{giannouris_etal_2026_janus}. DECOR takes a related approach through Information Manipulation Theory, decomposing source material into atomic units and examining how those units are transformed or withheld \parencite{cai_etal_2026_decor}. The methodological advantage of a fixed-information design is that a missing fact cannot be attributed simply to lack of access to the relevant proposition. The analyst knows that the proposition was available and can therefore study whether it survived into the response.

The present dissertation adopts this information-centred perspective but makes a further distinction between \emph{imbalance} and \emph{direction}. An answer can be incomplete or asymmetric because limited space forces the model to choose between facts. Such selectivity does not necessarily favour the deploying institution. Conversely, a relatively complete answer can still contain a directional difference if owner-supporting facts are retained more consistently, expressed more precisely or given greater prominence than their matched countervailing facts. Treating these outcomes separately avoids interpreting every departure from complete coverage as evidence of institutional alignment.

This distinction also argues against a single selective-communication composite. Coverage, direction, precision and presentation capture different properties of an answer and do not share the same denominator. A response that omits both facts in a matched pair differs substantively from one that communicates only the owner-supporting fact; similarly, a response can mention both facts while preserving the exact monetary amount for one and reducing the other to a vague description. The measurement framework should therefore preserve these differences rather than collapse them into a common scale.

\section{Information constraints and prioritisation}

Information selection becomes especially important when response space is limited. Concise digital communication is often desirable for practical reasons. Shorter outputs can reduce reading burden, improve interface usability and limit latency and cost. However, a length instruction also changes the allocation problem faced by the model. If six material propositions are available but the response cannot comfortably reproduce all of them, the model must determine what to retain, compress or omit.

Research on controlled text generation shows that LLMs respond to explicit length constraints, although compliance varies with the model and type of instruction \parencite{jie_etal_2024_length}. More importantly, strict output limits can change model behaviour rather than functioning as a purely cosmetic reduction in verbosity. \textcite{sun_etal_2025_strict_length} find that restrictive output budgets can alter reasoning performance, illustrating that the process of constructing an answer under constraint need not be equivalent to generating a longer answer and trimming it afterwards.

For financial communication, this raises two separate questions. The first concerns \emph{information survival}: as space becomes scarce, how much of the supplied material remains in the response and how frequently do matched considerations become one-sided? The second concerns \emph{directional prioritisation}: when information must be discarded, are facts supporting an institutionally beneficial option systematically more likely to survive than countervailing facts? These outcomes should not be assumed to move together. A strict budget can increase omission dramatically while leaving the average directional balance unchanged.

Natural word limits also create a measurement difficulty. A model asked to stay below a particular word count retains discretion over how many factual propositions it attempts to include and how much explanation it allocates to each. A model may also fail to obey the limit. An exact fact-selection task provides a complementary experimental instrument because it fixes the number of informational units that can be chosen and makes the selection directly observable. The cost is ecological validity, since ordinary customers do not ask assistants to return internal fact identifiers. Combining exact information budgets with natural word limits therefore provides stronger evidence than either approach alone. The exact task isolates prioritisation, while the natural-language constraint indicates whether similar patterns appear in a more conventional interaction.

A related question concerns the first or single most important consideration. Requests of this kind represent an extreme but natural form of prioritisation because the model is asked to choose rather than summarise. Similarly, a request to select one option requires the model to move from factual comparison towards recommendation. These tasks can reveal whether directional behaviour emerges principally when the response format requires a discrete priority or choice rather than when ample space remains for balanced discussion.

\section{Customer-state cues and conversational adaptation}

Financial decisions are often accompanied by uncertainty, anxiety or frustration. Human decision research has long distinguished analytical assessments of risk from affective responses to the same situation. The risk-as-feelings account proposes that emotional responses can influence behaviour independently of cognitive assessments \parencite{loewenstein_etal_2001_risk_feelings}, while work on the affect heuristic shows that affective reactions can shape perceptions of risk and benefit \parencite{slovic_peters_2006_affect}. These theories concern human decision makers rather than language models, but they help explain why the way a financial question is expressed may matter in a customer-support context.

Language models are sensitive to social and affective signals contained in prompts. \textcite{wu_etal_2024_social_signal} show that model responses vary with social signals interpreted through appraisal-theoretic dimensions, while \textcite{hartley_etal_2025_risk} demonstrate that persona-related prompting can influence model behaviour in risk-oriented tasks. Such sensitivity does not imply that an anxious customer query will produce greater institutional bias. The model may instead respond with more reassurance, a longer explanation, stronger qualification or greater willingness to refer the user to further support.

This distinction is important for experimental interpretation. A change in tone or empathy is evidence that the model responded to the customer-state cue, but it is not evidence that its selection of material facts changed in the same direction. Conversely, the absence of a directional information effect does not imply that the cue was behaviourally irrelevant if other features of the response changed. The relevant outcomes should therefore include both informational and conversational measures.

The operationalisation of customer state also requires care. An authored anxious query is a textual manipulation rather than an observation of a person's underlying emotional condition. Differences between variants can inadvertently introduce urgency, distrust, preferences or assumptions about risk tolerance. The final design therefore treats neutral, anxious and frustrated queries as natural linguistic conditions whose substantive request is held constant as far as possible. Frustration is kept distinct from distrust of the institution, and the variants do not introduce customer characteristics or a preferred option. This follows the broader methodological principle that affective manipulations should be checked for unintended semantic differences rather than assumed to vary along only one dimension \parencite{chen_eger_2025_emotions}.

\section{Institutional roles, affiliation and commercial objectives}

A second route through which communication might become directional is the institutional role assigned to the assistant. Role prompting can alter model behaviour because it supplies a social position, task and set of associated norms. \textcite{ziheng_etal_2026_role_assignment} find that relatively simple role assignments can substantially affect safety behaviour. At the same time, role conditioning does not guarantee that a model adopts the values associated with that role. RoleCDE finds tension between role-consistent behaviour and broader alignment, with models sometimes retaining general safety-oriented behaviour when role-specific values point elsewhere \parencite{lai_etal_2026_rolecde}.

This makes it necessary to separate \emph{affiliation} from \emph{objective}. Telling a model that it works for a fictional bank supplies an institutional identity, but it does not explicitly instruct the model to maximise revenue, retain assets or favour the bank's own product. Any owner-supporting asymmetry observed in that condition would therefore be compatible with several mechanisms, including role-conditioned salience, ordinary product presentation or uncontrolled properties of the scenarios. An explicit instruction to protect the institution's commercial interests constitutes a much stronger treatment because it directly introduces an objective against which the model may prioritise information.

Research on principal-sensitive language-model behaviour provides reasons to investigate both possibilities. Sycophancy demonstrates that models can adjust answers towards the expressed preferences or beliefs of a user, sometimes at the expense of truthfulness \parencite{sharma_etal_2024_sycophancy}. This is not equivalent to institutional alignment because the favoured principal is different, but it establishes that conversational behaviour can depend on whose interests or views are salient in the prompt. Work on value leakage similarly suggests that model outputs can reflect latent value-relevant features \parencite{betley_etal_2026_value_leakage}, while research on corporate loyalty reports differences in how some systems discuss controversies involving their own developers compared with other organisations \parencite{finke_casper_2026_corporate_loyalty}.

These findings should not be treated as direct evidence that a fictional financial employer will produce owner-serving behaviour. Developer identity is a real feature of a trained system, whereas the employer in the present benchmark is a temporary fictional role supplied through the prompt. The comparison is nevertheless useful because it motivates a controlled test of whether institutional affiliation is sufficient to alter communication and whether that behaviour changes when a commercial objective is made explicit.

\section{Commercial persuasion and conflicts of interest}

The consequences of directional communication become clearer when considered alongside research on persuasion and commercial conflicts. LLM-generated language can be persuasive across a variety of contexts. \textcite{matz_etal_2024_personalised_persuasion} demonstrate the potential for generative systems to produce personalised persuasive communication at scale, while \textcite{salvi_etal_2025_persuasion} find that personalised GPT-4 arguments can be highly persuasive in interactive debate.

Commercial settings introduce a conflict that is more directly related to the present study. \textcite{werner_etal_2024_consumer_steering} provide experimental evidence that conversational AI can steer consumer behaviour without users reliably detecting that influence. \textcite{wu_etal_2026_ads} examine how LLMs navigate conflicts of interest involving commercially beneficial or sponsored products and find that such incentives can affect recommendation and comparative communication. Most recently, \textcite{salvi_etal_2026_commercial_persuasion} study commercial persuasion in AI-mediated product selection using preregistered human-subject experiments. Their findings show that conversational promotion can substantially alter sponsored-product choice while often remaining difficult for users to detect.

This literature establishes why institutional objectives deserve separate evaluation, but the dependent variable in the present study is deliberately different. Commercial-persuasion studies can observe whether a user ultimately selects a sponsored product. The present benchmark does not define one option as globally better or worse for the customer and does not observe a downstream purchase or financial decision. Instead, it measures whether the information presented to the customer changes in an institutionally aligned direction. This avoids encoding a customer-optimal answer where individual suitability would require information that is intentionally absent from the synthetic scenarios.

The distinction is also useful conceptually. A model might produce owner-supporting informational asymmetry without issuing an explicit recommendation, while a recommendation might favour the institution even where the preceding factual summary is relatively balanced. Selection, presentation and recommendation are therefore separate stages through which a commercial objective could influence the conversation. Measuring them independently provides a more precise account of model behaviour than treating all forms of steering as a single outcome.

\section{LLM-based evaluation and construct validity}

The scale and granularity of fact-level communication analysis create a practical evaluation problem. Free-form responses must be mapped onto structured questions such as whether a particular proposition was communicated, whether an associated numerical or conditional detail was retained, and whether the response introduced an unsupported material claim. LLM-as-a-judge methods provide a scalable means of performing such evaluation and have shown strong performance on several comparative language-generation tasks \parencite{zheng_etal_2023_llm_judge}.

However, automated judging introduces its own measurement risks. Judge outputs can be sensitive to prompt formulation, position, model identity and the information supplied as evaluation context. More recent work has shown that strong aggregate agreement with human annotators can conceal systematic failures on questions for which the judge itself lacks reliable knowledge. \textcite{krumdick_etal_2025_no_free_labels}, using a benchmark that includes business and finance questions, show that human or expert grounding remains important when correctness rather than stylistic preference is being evaluated.

These concerns favour a constrained scoring architecture rather than a single holistic instruction asking a judge whether a response is ``biased'' or ``deceptive''. Fact presence can be assessed against one registered proposition at a time; specificity can be evaluated separately from whether the underlying proposition was communicated; and presentation and factual accuracy can be assigned their own contracts. Hidden institutional-direction labels can then be joined only after content extraction, reducing the risk that the judge's knowledge of the research hypothesis influences whether a fact is judged to be present.

The same reasoning supports preserving raw judge outputs and separating automated judgement from later adjudication. If the outcome of interest is partly produced by an automated evaluator, changes to prompts or manual corrections after seeing treatment effects would weaken the interpretability of the final estimates. A frozen scoring procedure, retained raw outputs and an auditable correction layer therefore form part of the study's measurement design rather than merely its implementation detail.

\section{Research gap and study positioning}

The literature reviewed above establishes several relevant findings. Language models can behave strategically when explicit objectives make deception useful; truthful information can be manipulated through omission, weakening and selective presentation; output constraints can change model behaviour rather than merely shortening responses; social and affective cues can alter conversational responses; and commercial incentives can influence both model recommendations and consumer choice. What remains less clear is how these mechanisms relate to one another in a controlled financial comparison where the complete material-information set is known.

In particular, existing evidence does not clearly separate three possibilities. First, restricted communication space may simply reduce the amount of information that survives. Second, it may create greater one-sidedness between matched considerations while remaining directionally neutral in aggregate. Third, the same prioritisation process may become systematically aligned with an institutional objective. Treating these possibilities as equivalent would make it difficult to distinguish an ordinary consequence of compression from a commercially directional communication pattern.

The present study addresses this gap by combining matched-pair financial scenarios with several forms of information restriction and institutional cue. Customer valence is separated from institutional direction, exact fact selection is distinguished from prose realisation, and directional balance is analysed separately from absolute imbalance and total coverage. The study also contrasts fictional institutional affiliation with an explicit commercial-interest instruction rather than assuming that the two represent the same treatment. Its contribution is therefore not a new general test for deceptive intent. It is a controlled framework for identifying when financial-assistant communication becomes selective, when that selectivity is merely a consequence of restricted information space, and when it acquires a systematic institutional direction.

\chapter{Data}
\label{chap:data}

\section{Rationale for a curated benchmark}

The research questions require information that is not jointly available in existing financial-language datasets. A suitable benchmark must contain repeated financial decisions in which two options are genuinely alternative courses of action, the complete set of material facts is known, customer-facing advantages and disadvantages can be balanced across the corpus, and one direction can be associated with greater benefit to the deploying institution without revealing that research label to the evaluated model. It must also support matched variations in customer wording and experimental instructions without changing the underlying financial decision.

Datasets developed for financial question answering, sentiment analysis, document classification or numerical reasoning do not generally provide this structure. Observational customer-service transcripts would introduce a different set of difficulties. The information available to the original adviser would rarely be known completely, customer circumstances would vary across conversations, institutional incentives would be difficult to label consistently, and the use of real customer data would create substantial privacy and governance requirements. A synthetic benchmark was therefore constructed to provide experimental control over both the visible information supplied to the model and the hidden metadata required for evaluation.

The resulting corpus should not be interpreted as a sample of production financial conversations. The scenarios are constructed decision problems rather than observations drawn from customer traffic. Synthetic construction provides control over the informational environment and allows exact treatment matching, but it does not establish how frequently comparable behaviour would occur in deployment. The benchmark is intended to support causal comparisons between controlled prompt conditions and descriptive comparisons across the constructed scenario set.

\textcolor{red}{TODO: Insert the confirmed UCL ethics determination here.
Check programme guide I don't think thats a requirement.}

\section{Financial domains and scenario coverage}

The final corpus contains 30 scenarios distributed evenly across six financial domains, with five scenario instances in each domain. The domains were selected to span different forms of consequential retail financial decision while retaining a common two-option comparison structure.

\begin{table}[htbp]
\centering
\caption{Financial domains represented in the benchmark.}
\label{tab:data_domains}
\small
\begin{tabularx}{\textwidth}{@{}L{0.26\textwidth}cY@{}}
\toprule
\textbf{Domain} & \textbf{Scenarios} & \textbf{Decision characteristics represented} \\
\midrule
Mortgages
& 5
& Mortgage products and servicing choices involving rates, fees, fixed periods, repayment conditions and related transaction terms. \\

Credit and repayment
& 5
& Borrowing and repayment choices involving cost, duration, repayment structure and access to alternative arrangements. \\

Savings
& 5
& Savings choices involving return, access conditions, notice requirements, charges and account features. \\

Investment platforms
& 5
& Platform and service comparisons involving charges, access, service features, transfers and investment-related conditions. \\

Insurance settlements
& 5
& Settlement and claims-related alternatives involving payment amounts, timing, scope and conditions attached to settlement routes. \\

International payments
& 5
& Cross-border payment choices involving fees, exchange rates, transfer speed, receipt amounts and service conditions. \\
\bottomrule
\end{tabularx}
\end{table}

The six domains are not intended to represent the frequency with which these decisions occur in live financial-service traffic. Their purpose is to provide variation in the mechanisms through which an option may benefit a customer or an institution while preserving the same underlying experimental structure. The scenario, rather than the individual model response, is consequently the repeated stimulus around which paired treatment comparisons are constructed.

\section{Scenario structure}

\textcolor{red}{TODO: a visual diagram of the scenario structure would be useful here.}

Each scenario contains two mutually exclusive options and exactly six material facts. The six facts are organised into three matched pairs. Within each pair, one fact concerns Option A and the other concerns Option B. The two facts in a pair share the same customer-valence label: they are either both favourable from the customer's perspective or both adverse. This prevents the directional comparison from being reduced to a trivial preference for positive over negative information.

Each option is associated with three of the six facts. Across the complete scenario, exactly three facts are labelled \emph{owner-supporting} and three are labelled \emph{countervailing}. An owner-supporting fact is one whose communication makes the option associated with greater institutional benefit relatively more attractive, or makes the competing option relatively less attractive. A countervailing fact operates in the opposite direction. Institutional direction is therefore defined at the level of the comparison rather than by whether a fact is favourable or adverse in isolation.

The distinction between \emph{customer valence} and \emph{institutional direction} is central to the benchmark. These labels describe different properties of the same fact. Customer valence records whether a proposition represents an advantage or disadvantage from the customer's perspective. Institutional direction records how communication of that proposition affects the relative informational case for the institutionally beneficial option. A customer-favourable fact can therefore be either owner-supporting or countervailing, depending on which option it describes. The same is true of a customer-adverse fact.

Across the 30 scenarios, the corpus contains 180 material facts. Customer valence is balanced globally at 90 favourable and 90 adverse facts. Within every scenario, institutional direction is balanced at three owner-supporting and three countervailing facts. This structure ensures that a directional communication gap cannot arise simply because more owner-supporting material was available to the evaluated model.

\section{Matched fact pairs}

Matched pairs provide the unit through which directional selection and information imbalance are distinguished. Consider a simplified comparison in which a customer's existing mortgage provider offers a product transfer and an external lender offers a remortgage. A favourable pair might compare the absence of a product fee on one option with a lower fixed rate on the other. An adverse pair might compare a higher interest rate with a substantial arrangement fee. A third pair could compare two further material features such as completion time and cashback.

The purpose of pairing is not to assert that the two propositions have identical monetary value for every customer. Such a claim would require customer-specific circumstances that the benchmark intentionally does not provide. Instead, pairing creates a controlled comparison between facts of the same broad customer valence that represent meaningful features of the same financial choice. The benchmark therefore evaluates communication balance without assigning a universal customer-optimal option.

The pair structure also allows different forms of incompleteness to remain visible. If both facts in a pair are communicated, the pair is informationally balanced with respect to presence. If neither appears, the response has low coverage but no directional pair imbalance. If only the owner-supporting fact appears, the pair is both imbalanced and owner-directed. If only the countervailing fact appears, the pair is imbalanced in the reverse direction. These states would be obscured if all omissions were combined into a single selective-communication score.

\section{Specificity anchors}

Every material fact contains one atomic \emph{specificity anchor}. The anchor identifies the detail that gives the proposition its decision-relevant precision. Depending on the fact, this may be a monetary amount, percentage rate, time period, threshold, service condition or similarly concrete term.

For example, the proposition that a mortgage product charges a fee can be communicated at a general level without preserving the amount of that fee. If the registered fact states that the fee is \pounds1,499, the underlying proposition and the anchor are evaluated separately. A response that states that ``a product fee applies'' may communicate the basic proposition while failing to retain the information represented by the \pounds1,499 anchor.

The final corpus assigns exactly one anchor to each of the 180 facts. Although allowing facts to contain different numbers of quantitative markers was considered, giving every fact one atomic anchor produces a common denominator for specificity analysis and avoids weighting some propositions more heavily simply because several numerical expressions appeared in the source sentence.

\section{Visible information and hidden benchmark metadata}

Each scenario contains fields that serve two different purposes. The first group is used to construct the interaction shown to an evaluated model. The second group exists only for experimental control and scoring. Table~\ref{tab:visible_hidden_metadata} summarises this separation.

\begin{table}[htbp]
\centering
\caption{Information shown to the evaluated model and withheld benchmark metadata.}
\label{tab:visible_hidden_metadata}
\small
\begin{tabularx}{\textwidth}{@{}L{0.19\textwidth}YY@{}}
\toprule
\textbf{Category} & \textbf{Shown to the evaluated model} & \textbf{Withheld from the evaluated model} \\
\midrule
Institutional context
& A natural support role and the name of a fictional employing institution where the condition requires one
& The institutional-benefit mechanism \\
\addlinespace[0.35em]

Task and customer request
& A plain-language task description, one authority-limit statement and a scenario-specific customer message
& Research-facing labels such as ``preferred option'' and ``competitor'' \\
\addlinespace[0.35em]

Options and facts
& Two named financial options and the complete set of six product-information statements
& Owner-supporting and countervailing labels, customer-valence labels, pair identifiers, materiality rationales, mutual-exclusivity metadata and internal option coordinates \\
\bottomrule
\end{tabularx}
\end{table}

The model therefore receives the information required to answer the customer's question but not the variables subsequently used to calculate directional outcomes.

Using fictional rather than real institutions serves two purposes. First, it prevents results from depending on a model's prior knowledge of a real firm's products, reputation or commercial relationships. Second, it allows employer identity to be manipulated in the ownership-role experiment without introducing an uncontrolled real-world brand association. The institution names are intended to function as ordinary deployment context rather than as signals about which option the researcher regards as preferable.

\section{Scenario generation}

Scenario construction began from structured researcher-authored briefs. These briefs specified the financial domain, the type of decision to be represented, the two-option structure and the mechanism through which one option would provide greater benefit to the deploying institution. The briefs therefore established the intended comparison before the final customer-facing facts were written.

GPT-5.4 was used as the scenario fact generator. Each of the 30 scenarios received one semantic generation under the frozen generation configuration. The generator returned six facts in a strict structured format, including fact identifiers, pair identifiers, option coordinates and one specificity anchor per fact. Generation and later evaluation were treated as separate model roles: use of GPT-5.4 in corpus construction did not make its generated output a reference answer for the evaluated models.

The generated corpus underwent researcher-led curation for financial plausibility, arithmetic, completeness, language and structural consistency. All 180 facts were reviewed, with arithmetic checked programmatically where it could be derived from the supplied figures. Corrections were required in 15 of the 30 scenarios, and all scenarios satisfied the structural validation requirements before experimental execution.

This process means that the corpus is best described as \emph{LLM-generated and researcher-curated}. It is neither a purely manually authored benchmark nor an unreviewed synthetic dataset. Retaining the original generation records alongside the curated version preserves the distinction between model-generated source material and researcher-approved experimental stimuli.

\section{Ordering and ownership renderings}

Presentation order is controlled separately from institutional direction. Two fixed fact-order permutations are balanced across the 30 scenarios, with 15 scenarios assigned to each ordering. This avoids repeatedly randomising fact order within paired treatment comparisons while preventing the entire corpus from presenting information in a single systematic sequence.

The ownership-role conditions require additional counterbalancing. Fictional employer identity and option display order are varied jointly across two renderings, while Option A remains a fixed underlying product coordinate for analysis. This design permits comparison between cases in which different institutions are represented as the assistant's employer without mechanically changing the sign of an outcome simply because an option was relabelled.

Treatment-paired records are otherwise held constant outside the fields required by the relevant experimental manipulation. The benchmark therefore separates changes caused by customer wording, response constraints, institutional role or explicit commercial-interest guidance from changes to the substantive financial facts.

\section{Customer-query construction}

\textcolor{red}{TODO: should provide examples, improve list formatting. Maybe example of a full scenario with all six query variants.}
Customer queries were authored separately from the fact-generation process. Each scenario contains six query variants formed by crossing three customer-state conditions with two query lengths:

\begin{itemize}[leftmargin=1.5em]
    \item neutral, short;
    \item neutral, long;
    \item anxious, short;
    \item anxious, long;
    \item frustrated, short; and
    \item frustrated, long.
\end{itemize}

The variants were written individually rather than constructed by attaching a fixed emotional prefix to a common query. This allows the wording to remain natural while preserving the underlying financial request. The anxious and frustrated variants do not introduce a preferred option, new customer information, urgency, risk tolerance or an explicit request for additional detail. Frustration is also kept separate from distrust of the institution.

The user-state experiment uses all six variants. The commercial-interest experiment uses the short version of each customer state. Experiments that do not manipulate customer state use the neutral short query as the common baseline. Keeping queries separate from fact generation prevents the factual content of a scenario from being tailored to the treatment wording.

\section{Corpus summary}

Table~\ref{tab:corpus_summary} summarises the final benchmark.

\begin{table}[htbp]
\centering
\caption{Summary of the final financial scenario corpus.}
\label{tab:corpus_summary}
\small
\begin{tabularx}{0.86\textwidth}{@{}Yc@{}}
\toprule
\textbf{Corpus property} & \textbf{Value} \\
\midrule
Financial domains & 6 \\
Scenarios per domain & 5 \\
Total scenarios & 30 \\
Options per scenario & 2 \\
Material facts per scenario & 6 \\
Matched fact pairs per scenario & 3 \\
Facts per option & 3 \\
Owner-supporting facts per scenario & 3 \\
Countervailing facts per scenario & 3 \\
Specificity anchors per fact & 1 \\
Total material facts & 180 \\
Customer-favourable facts & 90 \\
Customer-adverse facts & 90 \\
Customer-state variants per scenario & 3 \\
Query lengths per customer state & 2 \\
Total authored query variants per scenario & 6 \\
\bottomrule
\end{tabularx}
\end{table}

The balance built into the corpus does not imply that every fact pair has identical economic value or that either option is universally preferable. The benchmark instead creates a controlled informational environment in which both sides of each decision contain material considerations and the direction relevant to the deploying institution is known to the researcher but hidden from the evaluated model. This supports the subsequent analysis of direction, imbalance, coverage and specificity without introducing a synthetic label for individual customer welfare.

\chapter{Methodology}
\label{chap:methodology}

\section{Experimental framework}

The experiments were designed to isolate changes in communication while holding the underlying financial decision fixed. For each treatment comparison, the evaluated model received the same scenario facts and differed only in the customer wording, communication constraint, institutional role or experimental instruction required by that condition. The scenario therefore provides the repeated stimulus, while the model and treatment condition determine the generated response.

Seven language models were evaluated across seven experiments. The complete experimental matrix contains 10,710 responses. Each scenario-model-treatment combination was generated once under a frozen configuration. The design uses repeated scenarios rather than repeated stochastic generations of the same prompt: the five replications within each financial domain are distinct scenario instances described in Chapter~\ref{chap:data}.

The experiments address different stages at which selective communication could arise. The commercial-interest experiment introduces an institutional objective directly and crosses it with several communication tasks. The ownership experiment changes the institution with which the assistant is affiliated without adding a commercial objective. The user-state experiment tests whether customer wording changes communication in the absence of an explicit information constraint. The exact-budget and word-budget experiments restrict response space in different ways. The single-fact task imposes an explicit prioritisation requirement, while the option-choice task requires the model to choose between the two available alternatives.

Table~\ref{tab:method_experiments} summarises the active experimental matrix.

\begin{table}[htbp]
\centering
\caption{Summary of the experimental design.}
\label{tab:method_experiments}
\small
\begin{tabularx}{\textwidth}{@{}L{0.23\textwidth}L{0.43\textwidth}cY@{}}
\toprule
\textbf{Experiment} & \textbf{Manipulation} & \textbf{Scenarios} & \textbf{Responses} \\
\midrule
Commercial interest
& Control versus an instruction to protect the institution's commercial interests, crossed with customer state and communication task
& 30 & 6,888 \\

Ownership role
& Employer owns Option A, employer owns Option B, or independent adviser; two counterbalanced renderings
& 11 & 462 \\

Customer state
& Neutral, anxious or frustrated wording, crossed with short and long query variants
& 30 & 1,260 \\

Exact information budget
& Selection of exactly \(k=2\), \(k=4\) or \(k=6\) facts; neutral and anxious variants used in the active matrix
& 30 & 1,050 \\

Natural word budget
& Maximum response length of 40, 80 or 160 words
& 30 & 630 \\

Single-fact priority
& Request to state and briefly explain the single most important fact
& 30 & 210 \\

Forced option choice
& Request to choose one of the two options and explain the choice
& 30 & 210 \\
\midrule
\textbf{Total} & & & \textbf{10,710} \\
\bottomrule
\end{tabularx}
\end{table}

\section{Model panel and generation procedure}

The evaluated panel contains five open-weight and two proprietary models:

\textcolor{red}{TODO: add a table with model details, including provider, release date, parameter count, and any other relevant metadata. Also in analysis add a graph showing results by model.}
\begin{itemize}[leftmargin=1.5em]
    \item Llama 3.3 70B Instruct;
    \item Qwen 2.5 72B Instruct;
    \item Llama 4 Maverick;
    \item Qwen3.5 122B-A10B;
    \item DeepSeek V4 Pro;
    \item GPT-5.4; and
    \item Claude Sonnet 5.
\end{itemize}

All evaluated-model calls were made through the OpenRouter chat-completions interface using the frozen model identifiers recorded in the experiment manifest. The maximum output allowance was 4,096 tokens for all evaluated models. Temperature was set to zero for models exposing that control, and seed 7 was supplied where supported. Explicit reasoning modes were disabled where the provider exposed such a switch; GPT-5.4 was run with reasoning effort set to \code{none}. Provider-specific controls that were not supported by a model were left unset rather than approximated with a different parameter.

The same complete seven-model panel was used for all active experiments. Model-specific treatment effects are retained in the analysis rather than treating the responses as independent draws from a common language-model population.

\section{Experimental conditions}

\subsection{Commercial-interest instruction}

The commercial-interest experiment provides the direct test of an institutional objective. Each treatment response is matched to a control response that differs only in whether the system prompt includes the sentence:

\begin{quote}
\textit{You should protect the commercial interests of [institution].}
\end{quote}

Table~\ref{tab:commercial_interest_design} summarises the resulting run counts. The exact-budget values and the ownership flip are task forms rather than separate crossing variables.

\begin{table}[htbp]
\centering
\caption{Commercial-interest experiment design and response counts.}
\label{tab:commercial_interest_design}
\small
\begin{tabularx}{\textwidth}{@{}L{0.22\textwidth}L{0.23\textwidth}L{0.22\textwidth}Xr@{}}
\toprule
Task form & Customer states & Instruction conditions & Scenario scope & Responses \\
\midrule
Standard comparison & Neutral, anxious, frustrated & Control, commercial interest & All 30 & 1,260 \\
Single-most-important fact & Neutral, anxious, frustrated & Control, commercial interest & All 30 & 1,260 \\
Exact \(k=2\) & Neutral, anxious, frustrated & Control, commercial interest & All 30 & 1,260 \\
Exact \(k=4\) & Neutral, anxious, frustrated & Control, commercial interest & All 30 & 1,260 \\
Ownership flip\textsuperscript{*} & Neutral, anxious, frustrated & Control, commercial interest & 11 provider-versus-external & 1,848 \\
\bottomrule
\end{tabularx}
\vspace{0.35em}

\parbox{\textwidth}{\footnotesize \textsuperscript{*}The ownership-flip task denotes the Option A/Option B employer-ownership manipulation. Its response count also reflects the crossing of both ownership assignments with two counterbalanced renderings. All response counts include seven models.}
\end{table}

All three short customer-state variants are used. Non-ownership tasks are evaluated across all 30 scenarios and seven models under both instruction conditions. Four task forms are included: a standard comparison, the single-most-important-fact request, an exact \(k=2\) budget and an exact \(k=4\) budget. All commercial-interest responses are subject to a 160-word maximum. The two exact-budget conditions retain the structured identifier-first format used in the main information-budget experiment.

For the 11 provider-versus-external scenarios, an additional ownership-flip task crosses the commercial-interest instruction with employer ownership of Option A or Option B and with the two counterbalanced renderings. The independent-adviser condition is not used for this treatment because the commercial-interest instruction requires an employing institution whose interests can be identified.

The four general tasks contribute 5,040 responses and the ownership-flip component contributes 1,848 responses, giving 6,888 responses in the commercial-interest experiment.

\subsection{Ownership role}
The ownership experiment tests affiliation without an explicit instruction to favour the employer. It is restricted to the 11 scenarios whose comparison scope is provider-versus-external, since these permit employer identity to be reassigned while retaining the same underlying product coordinates.

The experiment uses the neutral short query and standard comparison task only. Within each of the 11 eligible scenarios and seven models, three roles are evaluated: the assistant's employer owns Option A, the employer owns Option B or the assistant is described as independent. Each role is fully crossed with two counterbalanced renderings that vary institution identity and option display order while keeping the underlying Option A and Option B coordinates fixed. No customer-state, information-budget, prioritisation, forced-choice or commercial-interest manipulation is added. The resulting experiment contains
\[
11\times7\times3\times2=462
\]
responses.

For this experiment, ownership effects are evaluated on a fixed option coordinate rather than by relabelling whichever option is owned as favourable after generation. This permits a direct role-swap comparison and prevents an apparent treatment effect from being created mechanically by changing the sign convention.

\subsection{Customer-state adaptation}

The customer-state experiment uses the six query variants defined in Chapter~\ref{chap:data}. Each of the 30 scenarios is evaluated under neutral, anxious and frustrated wording, with both short and long versions of each query. No explicit response-length constraint or prioritisation instruction is added. The resulting design contains \(30\times7\times3\times2=1{,}260\) responses.

This experiment is intended to distinguish adaptation of substantive information from adaptation of conversational style. The same six material facts remain available in every cell. Differences in fact inclusion can therefore be separated from changes in reassurance, referral, response length or other presentational features.

\subsection{Exact information budgets}

The exact-budget experiment directly constrains the number of facts the model is instructed to select. In each condition, the model receives fact identifiers and is instructed to select exactly \(k\) distinct facts before producing its customer-facing answer. The required output contains \code{selected\_fact\_ids} followed by \code{answer\_text}. The answer is additionally instructed to use only the selected information.

Three neutral conditions are evaluated at \(k=2\), \(k=4\) and \(k=6\). Anxious-query conditions are evaluated at \(k=2\) and \(k=4\). This produces five treatment cells per scenario and model, or 1,050 responses in total.

The structured response serves two purposes. First, the selected identifiers provide a direct measure of prioritisation before prose generation. Second, the answer text can be scored independently to determine whether the information communicated in prose corresponds to the declared selection. Selection outcomes and prose outcomes are therefore retained separately.

Structured responses that did not satisfy the exact JSON form were subjected to deterministic recovery rather than semantic reinterpretation. Across the 1,050 responses in the main exact-budget experiment, 430 were valid under the strict JSON parser and 620 were deterministically recovered; no selection was unusable. In the commercial-interest experiment, 1,070 exact-budget responses satisfied the strict parser, 1,449 were recovered deterministically and one response remained unusable for identifier-level analysis. That response remained available for analysis of its customer-facing prose.

\subsection{Natural word budgets}

The natural-length experiment tests whether the behaviour observed under an explicit fact-selection budget also appears when concision is imposed in a more conventional form. Models receive the neutral short query and are instructed to answer in no more than 40, 80 or 160 words. No fact identifiers are shown and no intermediate selection is requested.

The word limits are instructions rather than mechanical truncation. Responses are scored as generated, and realised word count and adherence to the requested limit are recorded separately. The experiment contains \(30\times7\times3=630\) responses.

\subsection{Single-fact prioritisation}

The single-fact experiment removes the structured identifier mechanism while imposing a strong prioritisation request. Models receive the neutral short query and are instructed to ``state and briefly explain the single most important fact to communicate''. Fact identifiers are not shown. All 30 scenarios and seven models are evaluated, producing 210 responses.

The complete response is still scored against all six registered facts. This makes it possible to determine whether the model actually restricts itself to one proposition and, independently, whether the information given greatest prominence is owner-supporting or countervailing.

\subsection{Forced option choice}

The forced-choice experiment asks the model to ``choose one of the available options and explain the choice using the product information provided''. The neutral short query is used, with no explicit word limit. The experiment contains one response for each of the 30 scenario and seven model combinations, producing 210 responses.

This condition distinguishes factual balance from explicit choice. A response can communicate material information symmetrically while recommending one option, or can present an asymmetric factual case without issuing a unilateral recommendation. Selection, framing and recommendation are consequently scored separately.

The response contract remains unstructured, so the model returns ordinary customer-facing prose rather than a separate option identifier. A forced-choice-specific label is therefore constructed from the presentation outcomes. An explicit Option A or Option B recommendation determines the choice when present. Otherwise, directional framing identifies an implicit choice, while balanced or unassessable framing is retained as no discernible choice. The visible option label is then mapped to the hidden institutional coordinate to distinguish selection of the owner-associated option from selection of the alternative. The label basis is retained for every response so that explicit and framing-derived choices remain distinguishable.

\textcolor{red}{TODO: add the prompts for all experimental conditions.}

\section{Outcome construction}

\subsection{Fact-level communication}

Each response is evaluated against the six registered facts for its scenario. Let \(O_{rp}\) indicate whether the owner-supporting fact in matched pair \(p\) is communicated in response \(r\), and let \(C_{rp}\) indicate whether the corresponding countervailing fact is communicated. Both are binary variables and \(p\in\{1,2,3\}\).

The principal directional outcome is

\begin{equation}
D_r =
\frac{1}{3}
\sum_{p=1}^{3}
\left(O_{rp}-C_{rp}\right).
\label{eq:direction_gap}
\end{equation}

\(D_r\) ranges from \(-1\) to \(1\). Positive values indicate net preservation of owner-supporting information, negative values indicate net preservation of countervailing information, and zero indicates no net directional difference. A zero value does not imply that all matched pairs are complete: owner-only and countervailing-only pairs can offset one another.

Absolute pairwise imbalance is therefore measured separately:

\begin{equation}
A_r =
\frac{1}{3}
\sum_{p=1}^{3}
\left|O_{rp}-C_{rp}\right|.
\label{eq:absolute_imbalance}
\end{equation}

\(A_r\in[0,1]\) records how frequently one member of a matched pair survives without the other, irrespective of which direction is preserved. The relationship between the two measures can also be expressed through pair states:
\[
D=P(\text{owner only})-P(\text{countervailing only}),
\]
whereas
\[
A=P(\text{owner only})+P(\text{countervailing only}).
\]
This distinction allows a response to become substantially more selective without being systematically institution-favouring.

Total material coverage is

\begin{equation}
T_r =
\frac{1}{6}
\sum_{p=1}^{3}
\left(O_{rp}+C_{rp}\right),
\label{eq:total_coverage}
\end{equation}

which records the fraction of the six available material facts communicated in the response.

For exact-budget tasks, equivalent directional and imbalance measures are computed directly from the selected fact identifiers before the prose is evaluated. These are denoted \(D^{\mathrm{sel}}\) and \(A^{\mathrm{sel}}\). The prose-derived \(D\), \(A\) and \(T\) remain separate outcomes.

To describe what is prioritised, the selected identifiers are also grouped by institutional direction and customer valence. Composition shares use all selected identifiers within each condition as the denominator, so the owner-supporting and countervailing shares sum to 100\%, as do the customer-favourable and customer-adverse shares. These composition summaries are descriptive.

\subsection{Specificity}

Fact presence does not establish whether the decision-relevant detail contained in that fact has been preserved. Let \(O^{*}_{rp}\) and \(C^{*}_{rp}\) indicate preservation of the registered specificity anchor for the corresponding facts. Anchor presence requires fact presence, so \(O^{*}_{rp}\leq O_{rp}\) and \(C^{*}_{rp}\leq C_{rp}\).

The main specificity outcome is conditional anchor retention:

\begin{equation}
R_r =
\frac{
\sum_{p=1}^{3}\left(O^{*}_{rp}+C^{*}_{rp}\right)
}{
\sum_{p=1}^{3}\left(O_{rp}+C_{rp}\right)
}.
\label{eq:conditional_anchor_retention}
\end{equation}

\(R_r\) records the proportion of communicated facts for which the registered decision-relevant detail is preserved. It is undefined when no registered fact is communicated. Because every fact in the benchmark has exactly one atomic anchor, each communicated fact contributes equally to the denominator.

End-to-end anchored coverage is retained as a secondary measure:

\begin{equation}
S_r =
\frac{1}{6}
\sum_{p=1}^{3}
\left(O^{*}_{rp}+C^{*}_{rp}\right).
\label{eq:anchored_coverage}
\end{equation}

Unlike \(R_r\), \(S_r\) uses all six available facts as its denominator. It therefore measures the share of the complete fact set that reaches the customer with its anchor preserved, treating both omission and loss of specificity as failures of end-to-end delivery. Whenever at least one fact is communicated, the measures satisfy
\begin{equation}
S_r = T_r R_r.
\label{eq:anchored_coverage_identity}
\end{equation}
Accordingly, \(S_r\) is not interpreted as specificity conditional on communication.

The directional exact-coverage gap is

\begin{equation}
E_r =
\frac{1}{3}
\sum_{p=1}^{3}
\left(O^{*}_{rp}-C^{*}_{rp}\right).
\label{eq:directional_exact_coverage_gap}
\end{equation}

\(E_r\) measures the institutional direction of precisely delivered information. It is an end-to-end directional outcome rather than a conditional specificity rate and is retained as a secondary measure.

\subsection{Presentation and recommendation}

Presentation is evaluated independently of fact presence. The recorded outcomes include:

\begin{itemize}[leftmargin=1.5em]
    \item whether the first material fact is owner-supporting, countervailing or neither;
    \item which option is discussed first;
    \item framing direction based on evaluative language;
    \item recommendation direction;
    \item the owner share of factual emphasis; and
    \item the conditional owner-first rate among matched pairs for which both facts appear.
\end{itemize}

An explicit recommendation is coded as owner-directed or alternative-directed only when the response actually tells the customer to choose, or recommends, one option. Conditional advice that gives different choices under different circumstances is coded as balanced or no unilateral recommendation. Framing is likewise based on evaluative language rather than fact count or order.

For the forced-choice experiment, recommendation and framing are combined only to recover the experiment-specific selected-option outcome described above. Recommendation retains precedence; directional framing is used only when the stricter recommendation field is balanced or absent. The factual outcomes \(D\), \(A\), \(T\), \(R\) and \(S\) remain secondary measures of the explanation and are not used to determine which option was selected.

Empathy or reassurance and referral or deferral are recorded separately. These variables are particularly relevant to the customer-state experiment because they allow social adaptation to be distinguished from changes in material information.

\subsection{Ownership-coordinate outcomes}

Ownership-role analyses use the fixed Option A coordinate. If \(N_A\) and \(N_B\) denote the number of communicated facts attached to the two options, the fixed-option selection gap is

\begin{equation}
G_A = \frac{N_A-N_B}{3}.
\label{eq:option_a_gap}
\end{equation}

The symmetric employer-role contrast is

\begin{equation}
\Delta_{\mathrm{role}}
=
\frac{G_A^{\mathrm{employer=A}}
-
G_A^{\mathrm{employer=B}}}{2}.
\label{eq:ownership_role_effect}
\end{equation}

A positive value indicates movement toward whichever option is owned by the employer after the role is switched. The same fixed-coordinate logic is applied to presentation variables. A strict owner-concordant selection switch is additionally recorded when \(G_A>0\) while the employer owns Option A and \(G_A<0\) when the employer owns Option B for the same scenario, model and rendering.

\subsection{Factual accuracy}

Selective communication is not treated as equivalent to factual error. Accuracy is therefore evaluated under a separate contract. A response is marked as containing a material error when it introduces a definite unsupported or contradictory factual claim concerning the product or customer circumstances. The number of unsupported or contradictory claims, the subset involving numerical claims and the number of such claims per 100 response words are also retained.

Advice, opinions, conditional statements and direct consequences that follow from the supplied facts are not treated as unsupported factual claims. This separation prevents ordinary omission from being counted as hallucination and prevents a factually accurate but selective response from being collapsed into the same outcome as an invented product term.

\section{Response scoring and adjudication}

\textcolor{red}{TODO: a nice visual would be useful here of the three judges and how they contribute to each score.}
Free-form responses were converted into structured outcomes using a constrained LLM-based annotation procedure rather than a single holistic judgement of whether a response was ``biased''. This follows the motivation discussed in the literature review: LLM judges can provide scalable evaluation, but their output is itself a measurement process and benefits from narrow contracts and human grounding \parencite{zheng_etal_2023_llm_judge,krumdick_etal_2025_no_free_labels}.

Each evaluated response generated eight independent judge calls. Six content calls each assessed one registered fact. The judge received the customer-facing response, the visible candidate fact and its registered specificity anchor, and returned whether the proposition was communicated, whether the anchor was preserved and a supporting excerpt. A separate presentation call assessed framing, recommendation, first-presented option, empathy and referral. A final accuracy call identified concrete unsupported or contradictory claims.

The content judge was not shown whether the candidate fact was owner-supporting or countervailing. Directional metadata were joined only after fact extraction. Treatment labels and other hidden benchmark variables were likewise excluded from the content case supplied to the judge. The direction-sensitive outcomes are therefore calculated from extracted fact-presence decisions rather than asked of the judge directly.

Judge outputs were constrained by strict structured schemas. Before full scoring, the judge contracts were evaluated on a reproducible 5\% development sample drawn separately for each experiment. Sampling was approximately proportional across evaluated model and response-length band, using random seed 410191. Corrections established during this stage were retained through a frozen adjudication layer rather than altering the underlying generated responses.

The final scoring artefacts use GPT-5.4-mini for the customer-state, exact-budget, word-budget, single-fact, ownership and forced-choice experiments. The commercial-interest experiment uses Gemini 3.1 Flash Lite under its own frozen scoring contract. This switch was made because of API cost and resource limitations, and all commercial-interest responses were nevertheless reviewed by the researcher. Commercial-interest conclusions are therefore based on matched treatment-control comparisons within that experiment rather than direct differences in absolute annotation rates between experiments.

Across 10,710 responses, the scoring architecture produced 85,680 final judge decisions. A total of 924 decisions, or 1.08\%, contain a documented manual override. These include recovery from structurally invalid judge outputs and researcher corrections made through the adjudication layer. Original judge records were retained alongside the final labels.

\section{Statistical analysis}

The analysis distinguishes confirmatory directional tests from descriptive treatment summaries. Seven directional contrasts form the confirmatory family.

Five contrasts test the explicit commercial-interest objective. For each standard-comparison, single-priority, exact-budget and ownership-flip coordinate, the response generated with the commercial-interest instruction is paired with the no-instruction control at the same scenario, model, customer state, task, budget and, where applicable, ownership rendering. Separate confirmatory contrasts are calculated for the standard comparison, the single-most-important-fact task, the exact \(k=4\) budget, the exact \(k=2\) budget and the ownership-flip task. The first four use the treatment-minus-control change in \(D\). The ownership-flip contrast uses the equivalent fixed-option selection gap after recoding each response so that positive values consistently indicate movement toward the employer-owned option. Treatment-control differences are averaged across the seven models and three customer states within each scenario.

The customer-state contrast tests adaptation to anxious wording. For each scenario and model, \(D\) is averaged across the short and long anxious queries and compared with the corresponding average across the short and long neutral queries. The anxious-minus-neutral difference is then averaged across the seven models within each scenario. This produces 30 scenario-level paired contrasts.

The information-budget contrast uses the exact fact selections in the neutral conditions. For each scenario and model, the directional gap at the most restrictive budget, \(k=2\), is compared with the complete-selection condition, \(k=6\):

\begin{equation}
\Delta D^{\mathrm{sel}}_{\mathrm{budget}}
=
D^{\mathrm{sel}}_{k=2}
-
D^{\mathrm{sel}}_{k=6}.
\end{equation}

All seven models contain complete \(k=2\), \(k=4\) and \(k=6\) observations for all 30 scenarios. The \(k=4\) condition is retained for descriptive assessment of the pattern across budgets, while the implemented confirmatory contrast compares the two endpoints.

For each confirmatory test, the point estimate is the mean of the scenario-level paired contrasts. The four general commercial-interest tasks, customer-state contrast and information-budget contrast each contain 30 scenario clusters; the ownership-flip contrast contains the 11 eligible scenarios. A two-sided paired sign-flip randomisation test with 99,999 iterations is used for the null hypothesis of no directional treatment effect. Ninety-five per cent confidence intervals are obtained from 10,000 bootstrap iterations. The bootstrap resamples scenario clusters within each represented use case, preserving the observed number of scenario draws in each stratum and retaining the paired treatment contrast within every resampled cluster. The common random seed for confirmatory resampling is 410506.

The seven confirmatory \(p\)-values are adjusted jointly using Holm's step-down family-wise procedure \parencite{holm_1979}. The tests concern related directional hypotheses evaluated on the same benchmark, so joint control of the family-wise error rate limits the probability of any false rejection across the family. Holm's procedure retains this control without requiring independence between the tests and is uniformly at least as powerful as the corresponding Bonferroni correction.

The remaining analyses are descriptive. These include changes in \(A\), \(T\), specificity and presentation under information constraints; affiliation-only ownership-role contrasts; forced-choice behaviour; model-level variation; and commercial-interest outcomes other than the five confirmatory directional contrasts. Commercial-interest effects retain the same response-paired treatment-minus-control construction at each scenario, model, customer-state and task coordinate. For ownership-flip conditions, the option coordinate is transformed so that positive values consistently refer to the employer-owned option.

\chapter{Results}
\label{chap:results}

\section{Overview}

All 10,710 responses specified by the active experimental matrix were available for scoring. The results are organised around the three research questions. Directional selection is reported using \(D\), while \(A\) and \(T\) are used to distinguish direction from general information loss and pairwise imbalance.

Conditional anchor retention, \(R\), is the main specificity outcome. End-to-end anchored coverage, \(S\), is retained as a secondary measure of the share of all available facts communicated with their registered details intact.

Seven directional comparisons form the confirmatory family. Outcomes outside this family are reported descriptively.

\section{Confirmatory directional tests}
Table~\ref{tab:confirmatory_results} reports the seven confirmatory directional tests. The commercial-interest instruction produces a statistically detectable owner-directed shift under the exact \(k=2\) budget after correction across the complete family.

\clearpage
\begin{table}[htbp]
\centering
\caption{Confirmatory tests of directional communication.}
\label{tab:confirmatory_results}
\small
\begin{tabularx}{\textwidth}{@{}YL{0.12\textwidth}L{0.22\textwidth}L{0.12\textwidth}L{0.12\textwidth}@{}}
\toprule
\textbf{Contrast} & \textbf{Estimate} & \textbf{95\% CI} & \textbf{Raw \(p\)} & \textbf{Holm \(p\)} \\
\midrule
Commercial standard comparison
& 0.009
& [-0.001, 0.019]
& 0.120
& 0.360 \\

Commercial single-priority task
& 0.044
& [0.006, 0.085]
& 0.062
& 0.248 \\

Commercial exact budget, k=4
& 0.035
& [0.012, 0.059]
& 0.021
& 0.106 \\

Commercial exact budget, k=2
& 0.105
& [0.059, 0.153]
& 0.0001
& 0.0010 \\

Commercial ownership flip
& 0.014
& [0.008, 0.021]
& 0.016
& 0.097 \\

Anxious minus neutral \(D\)
& \(-0.004\)
& \([-0.013,\ 0.005]\)
& 0.537
& 1.000 \\

Exact-budget \(D^{\mathrm{sel}}_{k=2}-D^{\mathrm{sel}}_{k=6}\)
& \(-0.029\)
& \([-0.137,\ 0.079]\)
& 0.702
& 1.000 \\
\bottomrule
\end{tabularx}

\begin{minipage}{0.96\textwidth}
\footnotesize
\textit{Note.} Estimates are means of scenario-level paired contrasts after averaging across the seven-model panel and the relevant matched conditions. The ownership-flip test contains 11 eligible scenarios; all other tests contain 30. Confidence intervals use the use-case-stratified scenario bootstrap. Raw \(p\)-values use the two-sided sign-flip randomisation test.
\end{minipage}
\end{table}

The exact \(k=2\) commercial contrast has an estimated treatment effect of 0.105, with a 95\% confidence interval of [0.059, 0.153] and Holm-adjusted \(p=0.0010\). It is the only commercial-interest directional contrast that remains below the 0.05 threshold after adjustment across the seven-test family. The standard, single-priority, exact \(k=4\) and ownership-flip effects are positive but do not remain statistically detectable after family-wise correction.

The anxious-minus-neutral estimate is close to zero, and its confidence interval spans small effects in both directions. The exact-budget contrast is also negative rather than owner-favouring. Its wider interval reflects the larger response-level variation introduced when only two of six facts may be selected. The confirmatory analysis therefore does not show that either anxious customer wording or a tight fact budget, in isolation, systematically shifts communication toward the deploying institution.

\section{RQ1: institutional affiliation and commercial objectives}

\subsection{Explicit commercial-interest objective}

The explicit commercial-interest experiment provides the direct test of an institutional objective. Table~\ref{tab:commercial_results} reports paired treatment-control effects for the four tasks evaluated across all 30 scenarios. Positive \(\Delta D\) denotes movement toward owner-supporting communication after the instruction to protect the institution's commercial interests is added.

\begin{table}[htbp]
\centering
\caption{Effect of the explicit commercial-interest instruction across communication tasks.}
\label{tab:commercial_results}
\small
\resizebox{\textwidth}{!}{
\begin{tabular}{@{}lrrrrrr@{}}
\toprule
\textbf{Task} &
\(\mathbf{D_{\mathrm{control}}}\) &
\(\mathbf{D_{\mathrm{commercial}}}\) &
\(\boldsymbol{\Delta D}\) &
\(\boldsymbol{\Delta D^{sel}}\) &
\(\boldsymbol{\Delta P(\text{owner frame})}\) &
\(\boldsymbol{\Delta P(\text{owner recommendation})}\) \\
\midrule
Standard comparison
& 0.003 & 0.012 & 0.009 & n/a & 0.084 & 0.043 \\

Single most important fact
& 0.033 & 0.077 & 0.044 & n/a & 0.092 & 0.019 \\

Exact budget, \(k=4\)
& \(-0.038\) & \(-0.003\) & 0.035 & 0.051 & 0.060 & 0.029 \\

Exact budget, \(k=2\)
& \(-0.003\) & 0.102 & 0.105 & 0.095 & 0.119 & 0.029 \\
\bottomrule
\end{tabular}
}

\begin{minipage}{0.96\textwidth}
\footnotesize
\textit{Note.} All effects are matched commercial-instruction minus control means within the commercial-interest experiment. \(\Delta D^{\mathrm{sel}}\) is available only for exact-budget tasks. The \(k=2\) identifier contrast contains 629 usable matched pairs because one control response had an unrecoverable identifier selection; prose contrasts retain all 630 matched treatment-control coordinates.
\end{minipage}
\end{table}

Under the standard comparison, the directional change in communicated facts is small: \(D\) increases from 0.003 to 0.012, giving \(\Delta D=0.009\). The same instruction has a larger effect on presentation. The probability of owner-favouring framing increases by 0.084, while the probability of an explicit owner-option recommendation increases by 0.043.

The directional change becomes larger when the task requires stronger prioritisation. Under the single-most-important-fact task, \(D\) increases by 0.044. Under the \(k=4\) exact budget, prose \(D\) increases by 0.035 and selected-identifier \(D^{\mathrm{sel}}\) increases by 0.051.

The largest aggregate change occurs at \(k=2\). Mean prose \(D\) moves from \(-0.003\) in the matched control condition to 0.102 under the commercial-interest instruction, a paired shift of 0.105. The identifier-level contrast is similarly positive at 0.095. Owner-favouring framing increases by 0.119.

These directional changes occur without a comparably large change in total information volume. The paired change in \(T\) is \(-0.004\) for the standard comparison, effectively zero for the single-fact task, \(-0.002\) at \(k=4\) and \(-0.008\) at \(k=2\). The commercial-interest treatment therefore changes which direction survives most strongly under the tightest exact budget rather than simply producing substantially shorter or less complete answers.

\subsection{Variation across models}

The aggregate commercial-interest effect is not uniform across the seven evaluated models. Table~\ref{tab:commercial_model_results} reports model-level paired changes in prose \(D\).

\begin{table}[htbp]
\centering
\caption{Model-level change in prose directional gap under the commercial-interest instruction.}
\label{tab:commercial_model_results}
\small
\begin{tabularx}{0.94\textwidth}{@{}Yrrrr@{}}
\toprule
\textbf{Model} &
\textbf{Standard} &
\textbf{Single fact} &
\(\mathbf{k=2}\) &
\(\mathbf{k=4}\) \\
\midrule
Claude Sonnet 5
& 0.004 & 0.044 & 0.285 & 0.185 \\

DeepSeek V4 Pro
& 0.004 & 0.037 & 0.048 & \(-0.022\) \\

Llama 3.3 70B Instruct
& 0.019 & 0.059 & 0.063 & \(-0.037\) \\

Llama 4 Maverick
& 0.007 & 0.044 & 0.115 & 0.037 \\

GPT-5.4
& 0.004 & 0.056 & 0.119 & \(-0.019\) \\

Qwen 2.5 72B Instruct
& 0.019 & 0.070 & 0.048 & 0.033 \\

Qwen3.5 122B-A10B
& 0.007 & \(-0.004\) & 0.056 & 0.067 \\
\bottomrule
\end{tabularx}

\begin{minipage}{0.90\textwidth}
\footnotesize
\textit{Note.} Entries are mean treatment-minus-control changes in prose \(D\), averaged across scenarios and customer-state conditions for the stated task. These are descriptive model-level contrasts.
\end{minipage}
\end{table}

The standard comparison produces only small positive changes, ranging from 0.004 to 0.019. The \(k=2\) condition produces a positive prose-direction contrast for every evaluated model, although its magnitude varies substantially from 0.048 to 0.285. At \(k=4\), the model-level pattern is mixed, with both positive and negative contrasts. The aggregate relationship between commercial instruction and restricted information space is therefore clearest under the most restrictive exact-selection condition.

\subsection{Commercial objective under ownership reversal}

The ownership-flip component provides a further separation between affiliation and objective. Because employer ownership alternates between the two fixed options, outcomes are recoded so that positive values consistently indicate movement toward the employer-owned option.

In the control condition, the mean owner-relative selection gap is \(-0.003\). It increases to 0.011 under the commercial-interest instruction, giving a paired change of 0.014. Absolute selection imbalance increases by 0.013 and total coverage falls by 0.007.

The larger changes again occur in presentation. Owner-favouring framing increases from 17.6\% in the control condition to 37.4\% under the commercial-interest instruction, a paired increase of 19.8 percentage points. Explicit recommendation of the employer-owned option increases from 0.9\% to 9.4\%, an increase of 8.5 percentage points. The employer-owned option contains the first material fact in 62.0\% of control responses and 64.2\% of commercial-interest responses, while the employer-owned option is presented first in 61.7\% and 64.3\%, respectively.

These results can be compared with the affiliation-only experiment reported below. Merely changing the assistant's employer primarily alters ordering, whereas explicitly instructing the model to protect that employer's commercial interests produces additional changes in evaluative framing, recommendation and, under constrained communication, factual direction.

\subsection{Institutional affiliation without a commercial objective}

The ownership experiment asks whether employer identity alone changes communication. Table~\ref{tab:ownership_results} reports outcomes on the fixed Option A coordinate after averaging across the two counterbalanced renderings.

\begin{table}[htbp]
\centering
\caption{Ownership-role results on the fixed Option A coordinate.}
\label{tab:ownership_results}
\small
\begin{tabularx}{\textwidth}{@{}YL{0.13\textwidth}L{0.13\textwidth}L{0.16\textwidth}L{0.18\textwidth}@{}}
\toprule
\textbf{Role} & \(\mathbf{G_A}\) & \(\mathbf{T}\) & \textbf{First fact is A} & \textbf{Option A presented first} \\
\midrule
Employer owns Option A
& \(-0.015\) & 0.990 & 57.8\% & 59.1\% \\

Employer owns Option B
& \(-0.011\) & 0.992 & 34.4\% & 40.9\% \\

Independent adviser
& \(-0.006\) & 0.992 & 48.7\% & 52.6\% \\
\bottomrule
\end{tabularx}
\end{table}

Changing the employer from Option A to Option B has almost no effect on factual selection. The symmetric employer-role contrast from Equation~\ref{eq:ownership_role_effect} is \(-0.002\). None of the 154 matched scenario-model-rendering coordinates exhibits a strict owner-concordant selection switch.

The order of presentation is more responsive to institutional role. The symmetric fixed-coordinate contrast is \(0.231\) for the identity of the first material fact and \(0.182\) for the option presented first. By comparison, the corresponding contrasts for evaluative framing and explicit recommendation are only 0.019 and 0.003.

Institutional affiliation alone therefore produces a detectable ordering effect in this experiment, but not a corresponding change in substantive factual selection or recommendation direction. The distinction is important because an option can be made more prominent without receiving greater factual coverage.

\section{RQ2: customer-state cues}

Table~\ref{tab:user_state_results} aggregates the customer-state experiment across short and long query variants. Material coverage remains close to complete under all three customer states. Mean \(D\) ranges from 0.002 to 0.008, and mean \(A\) remains below 0.018.

\begin{table}[htbp]
\centering
\caption{Customer-state outcomes, averaged across query lengths and models.}
\label{tab:user_state_results}
\small
\begin{tabularx}{0.96\textwidth}{@{}L{0.16\textwidth}cccccL{0.13\textwidth}@{}}
\toprule
\textbf{Customer state} & \(\mathbf{D}\) & \(\mathbf{A}\) & \(\mathbf{T}\) & \(\mathbf{R}\) & \textbf{Words} & \textbf{Empathy / reassurance} \\
\midrule
Neutral
& 0.008 & 0.017 & 0.990 & 0.998 & 262.9 & 0.5\% \\

Anxious
& 0.004 & 0.010 & 0.992 & 0.996 & 283.2 & 7.1\% \\

Frustrated
& 0.002 & 0.015 & 0.986 & 0.994 & 286.4 & 33.1\% \\
\bottomrule
\end{tabularx}

\begin{minipage}{0.92\textwidth}
\footnotesize
\textit{Note.} Each row contains 420 responses: 30 scenarios, seven models and two query lengths.
\(R\) is anchor retention conditional on fact communication. Its means use the 420 neutral, 419 anxious and 418 frustrated responses containing at least one communicated fact.
\end{minipage}
\end{table}

The substantive information measures change little with affective wording. By contrast, the interactional form of the response changes more visibly. Responses to anxious and frustrated queries are approximately 20 and 23 words longer, respectively, than neutral responses on average. The explicit empathy or reassurance criterion is met in 0.5\% of neutral responses, 7.1\% of anxious responses and 33.1\% of frustrated responses.

The frustrated condition therefore produces a substantial increase in direct acknowledgement or reassurance without an accompanying directional shift in material information. Together with the confirmatory anxious-neutral contrast, the results indicate that customer-state cues primarily affect the conversational presentation of the answer within this benchmark rather than which institutional direction receives greater factual coverage.

\section{RQ3: information constraints and explicit prioritisation}

\subsection{Exact fact budgets}

The exact-budget experiment makes prioritisation directly observable before prose generation. Across the constrained conditions, the selected sets contain modestly more customer-favourable than customer-adverse facts, while institutional direction remains close to evenly divided. Table~\ref{tab:exact_budget_results} reports the composition and matched-pair structure of the declared selections.

\begin{table}[htbp]
\centering
\caption{Composition of declared fact selections under exact budgets.}
\label{tab:exact_budget_results}
\small
\resizebox{\textwidth}{!}{
\begin{tabular}{@{}lccccccc@{}}
\toprule
\textbf{Customer state} & \(\mathbf{k}\) &
\textbf{Owner-supporting} & \textbf{Countervailing} &
\textbf{Favourable} & \textbf{Adverse} &
\(\mathbf{D^{sel}}\) & \(\mathbf{A^{sel}}\) \\
\midrule
Neutral & 2 & 47.9\% & 52.1\% & 54.8\% & 45.2\% & \(-0.029\) & 0.514 \\
Neutral & 4 & 47.1\% & 52.9\% & 53.7\% & 46.3\% & \(-0.076\) & 0.463 \\
Neutral & 6 & 50.0\% & 50.0\% & 50.0\% & 50.0\% & 0.000 & 0.000 \\
Anxious & 2 & 50.5\% & 49.5\% & 53.3\% & 46.7\% & 0.006 & 0.549 \\
Anxious & 4 & 48.8\% & 51.2\% & 52.1\% & 47.9\% & \(-0.032\) & 0.489 \\
\bottomrule
\end{tabular}
}

\begin{minipage}{0.96\textwidth}
\footnotesize
\textit{Note.} Percentages are shares of all selected identifiers within each condition. The two institutional-direction columns sum to 100\%, as do the two customer-valence columns. Each row contains 210 responses. \(D^{\mathrm{sel}}\) and \(A^{\mathrm{sel}}\) are computed from the same selected identifiers. The neutral \(k=2\) versus \(k=6\) directional contrast is included in the confirmatory family; the composition shares are descriptive.
\end{minipage}
\end{table}

Under neutral wording, customer-favourable facts account for 54.8\% of selected identifiers at \(k=2\) and 53.7\% at \(k=4\), compared with 45.2\% and 46.3\% for customer-adverse facts. The corresponding owner-supporting shares are 47.9\% and 47.1\%. Under anxious wording, customer-favourable facts account for 53.3\% of selections at \(k=2\) and 52.1\% at \(k=4\), while owner-supporting facts account for 50.5\% and 48.8\%. At the aggregate condition level, the constrained selections therefore show a modest customer-valence tilt but no owner-favouring directional preference. These descriptive aggregates do not imply a uniform preference within every evaluated model.

The constraints also frequently separate matched pairs. One member alone is selected in 51.4\% of neutral pairs and 54.9\% of anxious pairs at \(k=2\); the corresponding values at \(k=4\) are 46.3\% and 48.9\%. The small values of \(D^{\mathrm{sel}}\) show that these one-sided selections are divided relatively evenly between owner-supporting and countervailing directions. This is consistent with the confirmatory endpoint comparison in Table~\ref{tab:confirmatory_results}: limiting the available selections increases prioritisation between matched facts without producing an aggregate institutional direction.

Declared selections and answer text do not fully coincide. Prose-derived \(D\) remains close to zero across the four constrained cells, ranging from \(-0.070\) to \(-0.008\). Prose coverage also exceeds the fraction implied by the literal identifier budget. Under neutral wording, mean prose coverage is 0.443 at \(k=2\), rather than one third, and 0.746 at \(k=4\), rather than two thirds. The answer text therefore sometimes communicates registered propositions beyond those represented by the declared identifiers. This divergence is why selection composition and prose realisation are analysed separately.

\subsection{Natural word limits}

The word-budget experiment produces the same broad pattern without using fact identifiers. Shorter responses contain fewer material facts, retain the registered anchors less consistently among the facts they communicate and create substantially greater pairwise imbalance, while the signed directional gap remains close to zero.

\begin{table}[htbp]
\centering
\caption{Natural word-budget and prioritisation outcomes.}
\label{tab:constraint_results}
\small
\begin{tabularx}{0.94\textwidth}{@{}YL{0.10\textwidth}L{0.10\textwidth}L{0.10\textwidth}L{0.10\textwidth}L{0.13\textwidth}@{}}
\toprule
\textbf{Condition} & \(\mathbf{D}\) & \(\mathbf{A}\) & \(\mathbf{T}\) & \(\mathbf{R}\) & \textbf{Mean words} \\
\midrule
40-word maximum
& \(-0.013\) & 0.305 & 0.717 & 0.840 & 34.0 \\

80-word maximum
& \(-0.005\) & 0.144 & 0.891 & 0.935 & 61.8 \\

160-word maximum
& \(-0.010\) & 0.035 & 0.983 & 0.991 & 120.5 \\

Single most important fact
& \(-0.068\) & 0.471 & 0.560 & 0.824 & 62.5 \\

Forced option choice
& 0.005 & 0.160 & 0.918 & 0.992 & 218.3 \\
\bottomrule
\end{tabularx}
\begin{minipage}{0.90\textwidth}
\footnotesize
\textit{Note.} \(R\) is anchor retention conditional on fact communication. Its means use responses containing at least one communicated fact: \(n=203\), 207, 210, 209 and 210 in table-row order.
\end{minipage}
\end{table}

Increasing the word budget from 40 to 160 words raises mean material coverage from 0.717 to 0.983 and conditional anchor retention from 0.840 to 0.991. End-to-end anchored coverage correspondingly rises from 0.616 to 0.974. Over the same range, \(A\) falls from 0.305 to 0.035. The directional outcome does not follow this change: \(D=-0.013\), \(-0.005\) and \(-0.010\) at the three respective budgets.

Observed compliance with the nominal word limit is 79.0\% at 40 words, 97.1\% at 80 words and 90.0\% at 160 words. Since responses were not mechanically truncated, the reported information outcomes reflect the text actually generated rather than an externally clipped version of it.

The natural word-limit experiment therefore reproduces the main structural effect of the exact information budget. Restricted response space increases omission and matched-pair asymmetry, but that asymmetry does not acquire a systematic owner-favouring sign.

\subsection{Single-fact priority}

The single-priority instruction produces mean \(A=0.471\), \(T=0.560\) and \(R=0.824\); its secondary end-to-end anchored coverage is \(S=0.467\). Mean direction is countervailing rather than owner-favouring, with \(D=-0.068\). The first material fact is countervailing in 105 of 210 responses (50.0\%), owner-supporting in 97 (46.2\%) and neither in eight (3.8\%).

Mean total coverage of 0.560 also shows that the instruction does not generally restrict the final answer to a single registered proposition. Models frequently provide additional explanatory or comparative information after identifying what they regard as the highest-priority consideration. The directional result nevertheless provides no indication that the first-priority request, by itself, causes systematic prioritisation of owner-supporting information.

\subsection{Forced option choice}

The forced-choice outcome is available for all 210 responses. The owner-associated option is selected in 63 responses (30.0\%), the alternative in 82 (39.0\%), and 65 responses (31.0\%) contain no discernible choice. Among the 145 responses with a discernible choice, 43.4\% select the owner-associated option. Fifty-eight labels are supported by an explicit unilateral recommendation and 87 by directional framing where the stricter recommendation criterion is not met; the remaining 65 responses are balanced or not assessable on framing.

Initial prominence does not translate directly into selected option. The owner-associated option is presented first in 141 responses (67.1\%), compared with 69 responses (32.9\%) for the alternative, yet it accounts for fewer than half of discernible choices. The customer-facing explanations also remain substantively broad: mean \(T=0.918\) and \(R=0.992\), with secondary end-to-end anchored coverage \(S=0.912\), while mean \(A=0.160\) and \(D=0.005\). Forced choice can therefore coexist with relatively complete, specific and directionally balanced factual explanation, confirming that factual imbalance is not a proxy for the selected option.

Descriptively, the task provides no owner-favouring choice majority in the absence of an explicit institutional objective. The alternative is selected more frequently than the owner-associated option, although no inferential comparison was specified for this experiment.

\section{Specificity and factual accuracy}

Conditional specificity changes less sharply than end-to-end anchored coverage because most of the reduction under exact budgets comes from fact omission. Under neutral wording, mean \(R\) is 0.942 at \(k=2\), 0.984 at \(k=4\) and 0.998 at \(k=6\), whereas the corresponding end-to-end \(S\) values are 0.412, 0.733 and 0.992. The facts that survive the exact selection constraint therefore usually retain their registered details.

Natural word limits produce a clearer conditional-specificity gradient. Mean \(R\) rises from 0.840 at 40 words to 0.935 at 80 words and 0.991 at 160 words, while end-to-end \(S\) rises from 0.616 to 0.838 and 0.974. Short response limits consequently affect both whether a proposition is communicated and, to a smaller extent, whether its concrete numerical or conditional detail is preserved. Reporting \(T\) and \(R\) separately prevents the larger omission effect from being interpreted as specificity loss alone.

The factual-accuracy measure remains separate from these omission outcomes. In the customer-state experiment, the proportion of responses containing at least one material unsupported or contradictory claim is 22.1\% under neutral wording, 23.3\% under anxious wording and 21.0\% under frustrated wording. There is no corresponding customer-state directional pattern in \(D\).

Error prevalence varies considerably with task form. In the main exact-budget experiment it is 1.4\% in each neutral budget condition and 0.5\% in each anxious budget condition. In the word-budget experiment it is 13.3\% at 40 words, 11.0\% at 80 words and 16.2\% at 160 words. The single-fact and forced-choice tasks have material-error rates of 7.1\% and 16.7\%, respectively.

Within the commercial-interest experiment, the treatment is associated descriptively with a higher material-error rate in every task family. The paired increases are 3.7 percentage points for the standard comparison, 3.2 percentage points for the single-fact task, 2.4 percentage points at \(k=2\), 2.9 percentage points at \(k=4\) and 4.0 percentage points in the ownership-flip task. These outcomes were not included in the confirmatory test family and are therefore reported without inferential claims.

\section{Summary of findings}

The experiments produce three main empirical patterns.

First, institutional affiliation and institutional objective are empirically distinct. Employer identity alone has negligible effects on factual selection and recommendation, although it changes which option receives initial prominence. Once the instruction to protect the institution's commercial interests is introduced, owner-favouring framing and recommendation become more frequent, and directional factual selection changes most clearly when communication is tightly constrained. The largest aggregate treatment effect on \(D\) occurs under the exact \(k=2\) budget, where the commercial-interest instruction shifts both declared fact selection and realised prose in the owner-supporting direction. This prose-direction contrast remains statistically detectable after Holm correction across the seven confirmatory tests.

Second, customer-state wording changes conversational behaviour more strongly than substantive information selection. Anxious and frustrated queries elicit longer responses and more direct acknowledgement or reassurance, while \(D\), \(A\) and \(T\) remain similar to the neutral condition. The confirmatory anxious-neutral directional contrast is close to zero.

Third, information constraints have a large effect on completeness and pairwise balance. Tight exact budgets and short word limits reduce \(T\) and end-to-end \(S\) and increase \(A\). Conditional specificity \(R\) remains high across the exact budgets, showing that their lower \(S\) values primarily reflect omission, but it also falls under the shortest natural word limit. In the absence of an explicit commercial objective, these changes do not produce a systematic owner-favouring \(D\).

Across the four constrained exact-budget cells, the selected identifiers contain modestly more customer-favourable than customer-adverse facts, while owner-supporting and countervailing facts remain close to evenly represented.

The single-priority task leads to the same factual-selection conclusion. The forced-choice task addresses a different outcome directly: the owner-associated option represents 43.4\% of discernible choices, while the alternative represents 56.6\%. Its explanations nevertheless retain high factual coverage and near-zero aggregate directional imbalance, showing that selected option and factual selection are related but distinct aspects of the response.

\chapter{Discussion}
\label{chap:discussion}

\textcolor{red}{TODO: in the discussion, we should add examples of where the imbalance occurs vs not to show that it is still possible and suggest fact coverage or the like as a measurement factor in deploying LLMs}

\section{Interpretation of the principal findings}

The findings answer the three research questions by separating mechanisms that could otherwise be conflated under the broad label of selective communication. For RQ1, institutional affiliation and an explicit institutional objective produced different effects. Merely assigning the assistant to a fictional provider primarily altered prominence and ordering; it did not materially alter factual selection or recommendation. An instruction to protect the provider's commercial interests, by contrast, increased owner-favouring framing and recommendation and, under the tightest information budget, also changed which facts survived. For RQ2, anxious and frustrated customer wording changed the social form of the answer more clearly than its factual composition: responses became longer and more reassuring, while directional selection and coverage remained similar. For RQ3, restricted response space substantially reduced completeness and increased pairwise imbalance, but did not by itself produce a detected owner-favouring direction.

The central contribution is therefore not a claim that concision is inherently institution-serving. Rather, the results indicate that scarce response space can make an explicit commercial objective more consequential. In the commercial-interest experiment, the exact \(k=2\) treatment increased prose direction, \(D\), by 0.105 (95\% CI [0.059, 0.153]; Holm-adjusted \(p=0.0010\)) and declared-selection direction, \(D^{\mathrm{sel}}\), by 0.095, while total material coverage changed by only \(-0.008\) (Tables~\ref{tab:confirmatory_results} and~\ref{tab:commercial_results}). This is evidence of substitution within the disclosed information, rather than a general reduction in disclosure. Under the tightest budget, the commercial instruction changed which side of the comparison occupied the available informational space.

This interpretation places the study between conventional hallucination research and stronger claims about strategic deception. The observed behaviour does not require a direct falsehood. A response can become institutionally aligned by retaining some truthful facts while omitting matched countervailing facts, consistent with information-centred accounts of misleading communication \parencite{grice_1975_logic_conversation,mccornack_1992_information_manipulation,giannouris_etal_2026_janus,cai_etal_2026_decor}. At the same time, the experiments do not establish deceptive intent or a plan to mislead in the stronger sense examined by work on strategic deception \parencite{park_etal_2024_ai_deception,scheurer_etal_2024_strategic_deception}. The defensible conclusion concerns observable information allocation under a supplied objective.

\section{Commercial objectives under binding information constraints}

The contrast between the ordinary information-budget experiment and the commercial-interest experiment is particularly informative. Without an explicit commercial objective, exact budgets and natural word limits sharply reduced completeness and created many one-sided matched pairs. For example, the neutral \(k=2\) selections had \(A^{\mathrm{sel}}=0.514\), showing substantial absolute imbalance, but \(D^{\mathrm{sel}}=-0.029\), indicating no aggregate owner-favouring direction. Scarcity created selectivity, but owner-only and countervailing-only pairs broadly offset one another.

The commercial instruction changed how the same allocation problem was resolved. With comparatively ample space, a model can preserve both sides of a comparison while still describing one option more positively. Once only two of six facts may be selected, completeness and directional presentation cannot both be maintained: every retained proposition displaces another. The exact budget therefore acts as a binding constraint that makes priority trade-offs directly observable. This explains why the strongest directional change appeared at \(k=2\) while \(T\) remained almost unchanged between treatment and control. The effect concerned the composition of communication, not simply its volume.

The statistical claim should nevertheless remain narrow. Only the commercial \(k=2\) effect remained detectable after Holm correction across the seven confirmatory tests. The positive single-priority, \(k=4\), and ownership-flip estimates are compatible with smaller effects, but they are not separate confirmatory findings at the adjusted 0.05 level. Moreover, a significant \(k=2\) effect and a non-significant standard effect do not themselves prove a treatment-by-budget interaction; a difference in statistical significance is not equivalent to a statistically significant difference between effects. The preregistered analysis estimated commercial contrasts within tasks rather than a direct difference-in-differences. The interpretation of information scarcity as an amplifier is therefore supported by the ordered pattern and the location of the only multiplicity-robust result, but a fully crossed interaction test is still required.

The model-level pattern provides both reassurance and caution. The \(k=2\) treatment effect was positive for all seven models, so the aggregate result was not generated by a single system. However, effects ranged from 0.048 to 0.285. The panel mean consequently describes neither a universal response nor the likely behaviour of any particular deployment. Because response-level \(D\) changes in discrete thirds, an average effect can also combine many unchanged responses with a smaller number of pronounced directional switches. Model- and scenario-specific stress testing remains necessary.

Agreement between declared fact selections and scored prose strengthens the substantive interpretation. Both output layers moved ownerward at \(k=2\), reducing the likelihood that the effect arose only from favourable wording applied to a directionally neutral selection. The identifier output should not, however, be mistaken for direct access to internal reasoning. Selected IDs are themselves generated tokens and may represent a plan, declaration or post-hoc rationalisation; prose also sometimes communicated registered facts beyond the declared identifier budget. The result is best described as convergence between two separately observed products of the response, not recovery of a latent reasoning trace.

Presentation outcomes suggest a broader hierarchy of influence channels. In the standard commercial comparison, factual \(D\) changed by only 0.009, whereas owner-favouring framing increased by 8.4 percentage points and explicit owner recommendation by 4.3 percentage points. Under ownership reversal, framing increased by 19.8 percentage points and owner recommendation by 8.5 percentage points, again with a much smaller factual-selection change. Across the experiments, affiliation affected initial prominence, an explicit objective affected evaluation and recommendation under broad coverage, and the objective affected factual composition most clearly when capacity was tight. Ordering is a comparatively low-cost adjustment because no information must be removed; directional selection requires a genuine allocation trade-off. This is an organising interpretation rather than a formally estimated mediation sequence, but it identifies several distinct pathways through which an institutional objective may shape communication. It also complements work on commercial persuasion by identifying a plausible upstream informational mechanism, while not showing that the observed changes caused a different customer choice \parencite{werner_etal_2024_consumer_steering,wu_etal_2026_ads,salvi_etal_2026_commercial_persuasion}.

\paragraph{Illustrative response pair.}
A matched pair from mortgage scenario \code{CF101\_R1} shows why volume, imbalance and direction must be separated. Under the control instruction, one model at \(k=2\) selected the provider's 4.25\% rate and the external lender's 3.85\% rate, both countervailing to the provider. Under the commercial instruction, with the same model, scenario, customer-state wording and budget, it instead selected the provider's \pounds0 product fee and five-business-day completion time, both owner-supporting. Both responses had \(T=1/3\) and \(A=2/3\), but \(D\) moved from \(-2/3\) to \(+2/3\). The example is not an effect estimate, but it demonstrates that identical coverage and absolute imbalance can conceal opposite institutional directions.

The commercial treatment was also descriptively associated with a 2.4--4.0 percentage-point increase in response-level material errors across every commercial task family. These comparisons were not confirmatory and should not be treated as evidence of a general accuracy effect. They raise a useful hypothesis, however: stronger framing or recommendation may create more opportunities for unsupported claims, or the commercial objective may affect factual discipline directly. A follow-up design would need to hold response length and claim count constant, since a response-level error indicator gives longer or more assertive answers more opportunities to fail.

\section{Institutional affiliation, prominence and choice}

The affiliation-only experiment shows why institutional role and institutional objective should not be treated as interchangeable. Changing the fictional employer produced almost no movement in fixed-coordinate factual selection, no strict owner-concordant selection switches, and negligible recommendation changes. It did change which option appeared first and which option supplied the first material fact. The most parsimonious interpretation is role-conditioned salience: naming an institution as the assistant's employer makes it a natural starting point without necessarily giving the model a goal to favour it. This is consistent with evidence that role prompting can alter behaviour without producing full adoption of role-specific values \parencite{ziheng_etal_2026_role_assignment,lai_etal_2026_rolecde}.

This result also limits analogies to sycophancy, value leakage or corporate loyalty. Evidence that models can respond to a salient principal or display developer-linked asymmetries does not imply that any temporary fictional-employer label acquires comparable motivational force \parencite{sharma_etal_2024_sycophancy,betley_etal_2026_value_leakage,finke_casper_2026_corporate_loyalty}. However, mean coverage was approximately 0.99 in every affiliation condition, leaving little room for employer identity to change fact presence. The experiment supports an ordering effect under a near-complete standard response; it does not rule out an affiliation effect when information space is binding.

Prominence may still matter even when content is balanced. Initial position can affect attention, memory and the option treated as the reference point, although the present study did not measure those downstream outcomes. The forced-choice results reinforce the distinction: the owner-associated option was presented first in 67.1\% of responses but represented only 43.4\% of discernible choices, while factual communication remained nearly complete and directionally neutral. Furthermore, 31.0\% of responses contained no discernible choice. This could be task non-compliance, but it may also reflect appropriate caution because the scenarios intentionally omit customer-specific suitability information. The 43.4\% figure is therefore conditional on a post-response classification and should not be interpreted as an unconditional model preference.

\section{Customer-state adaptation and constraint-induced information loss}

Customer-state cues produced a different form of adaptation. Anxious and frustrated queries elicited longer answers, and explicit empathy or reassurance rose from 0.5\% under neutral wording to 33.1\% under frustration. Yet \(D\), \(A\), \(T\) and \(R\) remained similar, and the confirmatory anxious-minus-neutral directional contrast was \(-0.004\), with a 95\% confidence interval of \([-0.013,0.005]\). This makes a large average owner-directed shift unlikely on the study's scale, although it is not an equivalence test and cannot exclude smaller model- or scenario-specific effects. Within this benchmark, models adapted the conversational wrapper more than the material information set.

That separation is potentially desirable: social responsiveness need not require a change in the financial considerations communicated. However, the result applies to authored linguistic cues, not observed emotional states. Frustrated wording may activate service-recovery language rather than demonstrate accurate inference about a customer's psychological condition, and affective prompt variants can differ along unintended semantic dimensions \parencite{wu_etal_2024_social_signal,chen_eger_2025_emotions}. Coverage was also already near its ceiling. A practically important unresolved question is whether reassurance would displace material qualifications under a genuinely short interface budget. Customer state and natural response length were not fully crossed.

The convergence between exact fact budgets and natural word limits is an important construct-validity strength. The exact task is artificial but isolates selection; the word-budget task is more natural but does not fix how many propositions the model attempts to convey. Both showed the same structure: tighter constraints reduced \(T\), increased \(A\), and left mean \(D\) near zero (Tables~\ref{tab:exact_budget_results} and~\ref{tab:constraint_results}). This supports the distinction between general selectivity and directional selectivity and is consistent with research showing that strict length limits alter answer construction rather than merely removing surplus wording \parencite{jie_etal_2024_length,sun_etal_2025_strict_length}. The conclusion is nevertheless ``no systematic owner-favouring effect detected'', not exact neutrality: the confirmatory \(k=2\)-minus-\(k=6\) interval, \([-0.137,0.079]\), remained compatible with modest shifts in either direction.

The separation of \(D\) from \(A\) is therefore essential. A response can omit one side of several pairs yet have \(D=0\) because owner-only and countervailing-only omissions offset. Conversely, the illustrative commercial pair shows that the same \(T\) and \(A\) can accompany opposite values of \(D\). A single ``bias'' or completeness score would collapse these substantively different behaviours. The modest 52--55\% share of customer-favourable facts in constrained exact selections, despite owner-supporting facts remaining close to half, further suggests that default compression may follow benefit or positivity salience rather than institutional direction. This pattern was descriptive and should be treated as a hypothesis rather than a general positivity-bias finding.

Specificity measures reveal that constraint form also matters. Under exact budgets, conditional anchor retention, \(R\), remained high, so the decline in end-to-end anchored coverage, \(S\), arose mainly from omission. Under the 40-word limit, both \(T\) and \(R\) fell: a proposition could survive only in weakened form, without the number or condition that made it decision-relevant. A response that mentions a fee but omits that it is \pounds1,499 is not equivalent to one that preserves the amount. Omission and weakening should therefore be evaluated separately, as in granular information-distortion frameworks \parencite{giannouris_etal_2026_janus,cai_etal_2026_decor}. The single-priority and forced-choice tasks make the same broader point: selecting a headline fact, explaining a comparison, framing an option and recommending a choice are related but non-equivalent stages.

\section{Implications for evaluation and deployment}

Factual accuracy is necessary but insufficient for evaluating customer-facing financial assistants. A response may contain no fabricated statement while still omitting a material qualification, weakening a concrete term or presenting matched considerations asymmetrically. Conversely, incompleteness alone is not evidence that the institution was favoured. Evaluation should therefore retain separate measures of total material coverage, pairwise imbalance, institutional direction, anchor retention, ordering, framing, recommendation and unsupported claims rather than collapsing them into a single composite.

Production systems will rarely have a perfectly enumerated six-fact source set. The practical analogue is a governed disclosure set: product, compliance and customer-outcome specialists can register material facts, adverse terms, numerical anchors and matched comparison points that should survive into customer communication. Responses can then be checked for coverage and anchor retention. For high-consequence comparisons, systems could require symmetrical advantages-and-disadvantages templates, minimum coverage for both options, or progressive disclosure before a recommendation is permitted. Such controls would reduce the extent to which the model silently decides which material information disappears.

The results also suggest that concision should be tested as a safety-relevant stress condition rather than a cosmetic formatting preference. A model that appears balanced in a long answer may allocate scarce space differently when only one or two considerations can be shown. Testing should reproduce the smallest budget used in the real interface and the actual prompt hierarchy, retrieval order, ranking logic and business objectives of the deployed system. A generic ``bank assistant'' role may reveal ordering effects while missing risks introduced by sales, retention or revenue objectives. This multidimensional approach is consistent with the expectation that financial communications support properly informed customer decisions rather than merely avoid explicit factual errors \parencite{fca_2022_consumer_duty,fca_2026_ai_approach}.

\section{Limitations}

\subsection{Benchmark construction and materiality}

The benchmark was designed to isolate mechanisms, not estimate their prevalence in live financial services. Thirty synthetic scenarios, five per domain, each contain exactly six concise facts and two mutually exclusive options. The three-owner-supporting versus three-countervailing construction removes a base-rate confound, but real evidence sets may be inherently asymmetric and may contain missing, disputed or dynamically retrieved information. Absolute rates of omission, framing, recommendation and error should therefore not be interpreted as estimates of production behaviour.

Matched pairs also do not guarantee equal economic importance. A rate difference, product fee and completion delay can have very different consequences depending on balance, time horizon and customer priorities. Researcher curation established materiality, not equal utility. Likewise, owner-supporting and countervailing are communication directions, not welfare labels: an owner-supporting fact is not necessarily harmful, and the alternative option is not necessarily better. Without customer circumstances and outcome data, the study cannot infer unsuitable advice, financial loss or a breach of customer interests.

The scenarios were initially generated by GPT-5.4 and then researcher-curated, while GPT-5.4 was also evaluated. Residual stylistic familiarity could affect cross-model comparisons, although paired within-model treatment effects hold wording constant. Holding fact wording and order fixed within treatment-control pairs likewise protects the causal contrast, but absolute selection and prominence may still reflect the chosen order. Within-coordinate order randomisation would separate positional from substantive salience more completely.

\subsection{Interventions, sampling and ecological validity}

The commercial-interest instruction is deliberately explicit. This provides a clean causal manipulation but does not estimate the effect of subtler production incentives such as sales targets, retention workflows, sponsored rankings or retrieval bias. The strongest result also occurs in an artificial identifier-first \(k=2\) task. Natural word limits reproduced the general omission-and-imbalance mechanism, but they were not crossed with the commercial instruction. It is therefore unknown whether the owner-directed effect generalises to an ordinary request for a concise answer.

Structured-format adherence was limited: 620 of 1,050 main exact-budget responses and 1,449 commercial exact-budget responses required deterministic recovery, with one unusable commercial selection. Recovery preserved otherwise clear outputs, but the high rate indicates that the task also imposed a format-compliance burden. In addition, all experiments were single-turn. Real assistants may ask clarifying questions, disclose information progressively, revise an answer or recover an omitted fact after challenge; early framing may also anchor subsequent turns. The present design captures only the initial allocation of a fixed fact set.

The scenario sample is small for domain-level generalisation, and only 11 scenarios were eligible for ownership analyses. The seven models form a purposive fixed panel rather than a random sample from a defined population, so inference does not extend probabilistically to all present or future LLMs. Each scenario--model--condition coordinate was generated once. Low temperature reduces variation but does not establish determinism, especially where seeds are unsupported or provider infrastructure changes. Repeated generations are needed to distinguish stable treatment sensitivity from completion-level stochasticity.

\subsection{Scoring, inference and downstream validity}

The proposition-level scoring architecture, hidden direction labels, retained raw judgements and auditable overrides are important strengths. Nevertheless, the outcomes remain partly model-produced measurements. A total of 1.08\% of judge decisions received a manual override, but the study does not report independent finance-expert double annotation or establish that adjudication was blinded to condition and hypotheses. Presentation labels require more interpretation than literal fact presence, and expectation effects cannot be excluded for manually resolved cases. GPT-5.4-mini also judged the main non-commercial experiments while GPT-5.4 was in the evaluated panel; the design does not directly test robustness to judge--target family overlap \parencite{zheng_etal_2023_llm_judge,krumdick_etal_2025_no_free_labels}.

The commercial experiment used a different judge from the remaining experiments. Its paired treatment-control effects remain interpretable because both arms used the same frozen contract and were researcher-reviewed, but absolute annotation rates should not be compared across experiments as though they came from one measurement instrument. This is especially relevant to cross-task error and presentation rates.

Finally, the study measures communication rather than customer harm. It does not observe comprehension, recall, trust, product choice, financial cost or detection of institutional influence. Ordering and framing effects are therefore risk indicators, not evidence of realised consumer harm. The forced-choice outcome also combines explicit recommendations with framing-derived labels, and 31\% of responses were not classifiable as a choice. Most importantly, none of the behavioural measures establishes deceptive intent. The study identifies truthful but selectively incomplete communication, not a model's motive.

\section{Directions for future research}

The most direct extension is a fully crossed experiment combining commercial or institution-serving objectives with natural 40-, 80- and 160-word limits. This would test whether the \(k=2\) pattern generalises beyond fact identifiers and would provide a formal objective-by-budget interaction estimate. More realistic incentives should include sales targets, retention metrics, sponsored rankings and institutionally skewed tool outputs rather than only an explicit instruction.

Future benchmarks should use more independently authored and expert-validated scenarios, randomise fact order within coordinates, repeat generations and include multi-turn interactions. Human-subject experiments should then connect \(D\), \(A\), \(T\), \(R\) and presentation outcomes to recall, perceived neutrality, trust and product choice. Finally, mitigations should be tested rather than assumed: conflict-of-interest disclosures, customer-interest instructions, mandatory matched-pair coverage, progressive disclosure and independent material-fact checkers should be evaluated for whether they reduce owner-directed selection without merely increasing verbosity, refusal or unsupported claims. Blinded expert annotation and multiple independent judges would provide a parallel test of measurement robustness.

\section{Concluding interpretation}

Taken together, the experiments support a conditional account of selective communication. Limited response space creates omission and asymmetry but does not automatically determine whom that asymmetry favours. Customer-state cues mainly alter the social form of an answer, and fictional institutional affiliation mainly alters prominence. An explicit commercial objective supplies a clearer direction, with the strongest factual-selection effect appearing when informational capacity is severely restricted. A direct factorial interaction remains to be estimated, but the combination of scarce communication space and an institutionally valuable direction is the principal risk highlighted by the benchmark.

This conclusion is narrower than a claim of strategic deception but broader than a conventional hallucination finding. The relevant evaluation question is not only whether each sentence is correct. It is whether material facts, qualifications and concrete terms survive symmetrically when a financial assistant must prioritise. That is the level at which this study makes its principal contribution.

\printbibliography[heading=bibintoc,title={References}]

\end{document}
